\chapter{The Thermal Instability of a Layer of Fluid Heated from Below}
\label{ch:3}

\section*{2. THE EFFECT OF ROTATION}

\section{Introduction}\label{sec:3-19}
In this chapter we shall investigate the effect of rotation on the simple problem of thermal instability considered in the last chapter. It will appear that rotation introduces a number of new elements into the problem; and some of its consequences are, at first sight, unexpected: the role of viscosity is, for example, inverted. The origin of this and other consequences of rotation can be traced to certain general theorems, relating to vorticity, in the dynamics of rotating fluids. For this reason, we shall preface the investigation of the stability by a discussion of these general theorems.

\section{The theorems of Helmholtz and Kelvin}\label{sec:3-20}
Consider an incompressible, inviscid, fluid. The equations governing it are:\footnote{The principal theorems to be established are not restricted to incompressible fluids; they apply to compressible fluids, also, if $p$ is a function of $\rho$ only. In the latter case $(\operatorname{grad} p)/\rho$ can be written as $\operatorname{grad} \chi$ where $\chi = \int dp/\rho$; and the quantity on the left-hand side of equation~\eqref{eq:3-1} continues to be the gradient of a scalar function. It is on this last circumstance that the principal theorems to be established depend; we shall not, in particular, have occasion to use the equation of continuity~\eqref{eq:3-2}.}
\begin{equation}\label{eq:3-1}
\frac{\partial u_i}{\partial t} + u_j \frac{\partial u_i}{\partial x_j} = -\frac{\partial}{\partial x_i}\!\left(\frac{p}{\rho} + V\right),
\end{equation}
and\dag
\begin{equation}\label{eq:3-2}
\frac{\partial u_i}{\partial x_i} = 0,
\end{equation}
where it has been assumed that the external forces are derivable from a potential $V$. The vorticity $\boldsymbol{\omega}$ is, of course, defined by
\begin{equation}\label{eq:3-3}
\boldsymbol{\omega} = \operatorname{curl} \mathbf{u} \quad \text{or} \quad \omega_i = \epsilon_{ijk}\frac{\partial u_k}{\partial x_j}.
\end{equation}

A curve drawn from point to point in the fluid, so that its direction is everywhere that of the instantaneous direction of $\boldsymbol{\omega}$, is called a \emph{vortex line}. The differential equations determining it are
\begin{equation}\label{eq:3-4}
\frac{dx_1}{\omega_1} = \frac{dx_2}{\omega_2} = \frac{dx_3}{\omega_3}.
\end{equation}

If through every point of a small closed curve we draw the corresponding vortex line, we obtain a tube which is called a \emph{vortex tube}.

Since
\begin{equation}\label{eq:3-5}
\operatorname{div} \boldsymbol{\omega} = 0,
\end{equation}
it follows from Gauss's theorem that the integral of the normal component of $\boldsymbol{\omega}$ over any closed surface is zero:
\begin{equation}\label{eq:3-6}
\oint_S \boldsymbol{\omega} \cdot d\mathbf{S} = 0 \quad (S\ \text{any closed surface}).
\end{equation}

Apply this result to the closed surface bounding an element of fluid contained in a section of the vortex tube (see Fig.~18). The sides of the vortex tube clearly do not contribute to the surface integral and equation~\eqref{eq:3-6}, in this case, gives
\begin{equation}\label{eq:3-7}
\int_{S_1} \boldsymbol{\omega} \cdot d\mathbf{S}_1 = \int_{S_2} \boldsymbol{\omega} \cdot d\mathbf{S}_2.
\end{equation}

\figurePlaceholder{18}{A vortex tube.}

\emph{The flux of vorticity across any section of a vortex tube is the same}; it, therefore, represents a characteristic of the tube. If we consider a vortex tube of infinitesimal cross-section, the flux across a section $d\mathbf{S}$ normal to itself is $|\boldsymbol{\omega}|dS$. By the theorem we have just proved, $|\boldsymbol{\omega}|dS$ is constant along the tube; it cannot therefore terminate at any point in the fluid. \emph{Vortex lines must either be closed, or terminate on the boundaries.}

We shall now obtain an equation of motion for the vorticity. First, we observe that
\begin{equation}\label{eq:3-8}
\epsilon_{ijk}\, u_j\, \omega_k = \epsilon_{ijk}\,\epsilon_{klm}\, u_j \frac{\partial u_m}{\partial x_l}
= (\delta_{il}\,\delta_{jm} - \delta_{im}\,\delta_{jl})\, u_j \frac{\partial u_m}{\partial x_l}
= u_j\frac{\partial u_j}{\partial x_i} - u_j\frac{\partial u_i}{\partial x_j},
\end{equation}
or
\begin{equation}\label{eq:3-9}
u_j\frac{\partial u_i}{\partial x_j} = \frac{1}{2}\frac{\partial}{\partial x_i}|\mathbf{u}|^2 - \epsilon_{ijk}\, u_j\, \omega_k.
\end{equation}

Using this relation in~\eqref{eq:3-1}, we have
\begin{equation}\label{eq:3-10}
\frac{\partial u_i}{\partial t} - \epsilon_{ijk}\, u_j\, \omega_k = -\frac{\partial}{\partial x_i}\!\left(\frac{p}{\rho} + \tfrac{1}{2}|\mathbf{u}|^2 + V\right).
\end{equation}

Letting
\begin{equation}\label{eq:3-11}
\varpi = p/\rho + \tfrac{1}{2}|\mathbf{u}|^2 + V,
\end{equation}
we can write
\begin{equation}\label{eq:3-12}
\frac{\partial \mathbf{u}}{\partial t} - \mathbf{u}\times\boldsymbol{\omega} = \operatorname{grad}\varpi.
\end{equation}

Taking the curl of this equation, we obtain
\begin{equation}\label{eq:3-13}
\frac{\partial\boldsymbol{\omega}}{\partial t} - \operatorname{curl}(\mathbf{u}\times\boldsymbol{\omega}) = 0;
\end{equation}
this is the required equation of motion for $\boldsymbol{\omega}$. The principal theorem on vorticity follows from this equation.

Consider a surface $S$ enclosed by a simple closed contour $C$. Let $d\mathbf{S}$ be an element of this surface. Multiplying equation~\eqref{eq:3-13} scalarly by $d\mathbf{S}$, and integrating over $S$, we obtain
\begin{equation}\label{eq:3-14}
\int_S \frac{\partial\boldsymbol{\omega}}{\partial t}\cdot d\mathbf{S} - \int_S \operatorname{curl}(\mathbf{u}\times\boldsymbol{\omega})\cdot d\mathbf{S} = 0.
\end{equation}

\figurePlaceholder{19}{Illustrating the Helmholtz--Kelvin theorem.}

Transforming the second of these integrals by Stokes's theorem, we have
\begin{equation}\label{eq:3-15}
\int_S \frac{\partial\boldsymbol{\omega}}{\partial t}\cdot d\mathbf{S} + \oint_C \boldsymbol{\omega}\cdot(\mathbf{u}\times d\mathbf{s}) = 0,
\end{equation}
where $d\mathbf{s}$ is an element of arc of the contour $C$. The physical meaning of this last equation becomes clearer if we multiply the equation by an infinitesimal interval of time $\Delta t$, and then interpret the different terms. We have
\begin{equation}\label{eq:3-16}
\Delta t \int_S \frac{\partial\boldsymbol{\omega}_t}{\partial t}\cdot d\mathbf{S} + \oint_{C_t} \boldsymbol{\omega}_t\cdot(\mathbf{u}\,\Delta t \times d\mathbf{s}) = 0,
\end{equation}
where we have inserted a subscript $t$ to the various quantities to emphasize that they are the instantaneous values at the time indicated (see Fig.~19). Now $\mathbf{u}\,\Delta t \times d\mathbf{s}$ is the area swept out by the fluid elements along $d\mathbf{s}$ in a time $\Delta t$. Accordingly, the integral along $C_t$ in~\eqref{eq:3-16} is, in reality, a \emph{surface integral} over the elementary strip connecting the contour $C_t$ with the contour $C_{t+\Delta t}$ to which $C_t$ will be carried in the time $\Delta t$. Applying the result~\eqref{eq:3-6} to the closed surface consisting of $S_t$ (bounded by $C_t$), $S_{t+\Delta t}$ (bounded by $C_{t+\Delta t}$), and the strip connecting $C_t$ and $C_{t+\Delta t}$, we have\footnote{Note that the validity of this equation depends on $\operatorname{div}\boldsymbol{\omega} = 0$.}
\begin{equation}\label{eq:3-17}
\int_{S_{t+\Delta t}} \boldsymbol{\omega}_t\cdot(\mathbf{u}\,\Delta t \times d\mathbf{s}) = \int_{S_{t+\Delta t}} \boldsymbol{\omega}_t\cdot d\mathbf{S} - \int_{S_t} \boldsymbol{\omega}_t\cdot d\mathbf{S}.
\end{equation}

Using this in equation~\eqref{eq:3-16}, we obtain
\begin{equation}\label{eq:3-18}
\Delta t \int_{S_t} \frac{\partial\boldsymbol{\omega}_t}{\partial t}\cdot d\mathbf{S} + \int_{S_{t+\Delta t}} \boldsymbol{\omega}_t\cdot d\mathbf{S} = \int_{S_t} \boldsymbol{\omega}_t\cdot d\mathbf{S}.
\end{equation}

With an error of order $(\Delta t)^2$, we can replace the integral over $S_t$ on the left-hand side by an integral over $S_{t+\Delta t}$; we thus find,
\begin{equation}\label{eq:3-19}
\int_{S_{t+\Delta t}} \!\left(\boldsymbol{\omega}_t + \Delta t\frac{\partial\boldsymbol{\omega}_t}{\partial t}\right)\!\cdot d\mathbf{S} = \int_{S_t} \boldsymbol{\omega}_t\cdot d\mathbf{S} + O(\Delta t^2),
\end{equation}
or
\begin{equation}\label{eq:3-20}
\int_{S_{t+\Delta t}} \boldsymbol{\omega}_{t+\Delta t}\cdot d\mathbf{S} = \int_{S_t} \boldsymbol{\omega}_t\cdot d\mathbf{S} + O(\Delta t^2).
\end{equation}

Passing to the limit $\Delta t = 0$, we obtain
\begin{equation}\label{eq:3-21}
\frac{d}{dt}\int_{S_t}\boldsymbol{\omega}\cdot d\mathbf{S} = 0.
\end{equation}

We have thus proved: \emph{The integral of the normal component of\/ $\boldsymbol{\omega}$ over any surface $S$ bounded by a closed curve remains constant as we follow the surface $S$ with the motion of the fluid elements constituting it}:
\begin{equation}\label{eq:3-22}
\int_{S_t}\boldsymbol{\omega}_t\cdot d\mathbf{S} = \text{constant}.
\end{equation}

This is the principal theorem on vorticity due to Helmholtz and Kelvin. We may state it also in the form: \emph{the strength of a vortex tube is an integral of the equations of motion}.

Transforming~\eqref{eq:3-22} by Stokes's theorem, we have
\begin{equation}\label{eq:3-23}
\int_S \boldsymbol{\omega}\cdot d\mathbf{S} = \int_S \operatorname{curl}\mathbf{u}\cdot d\mathbf{S} \stackrel{*}{=} \oint_C \mathbf{u}\cdot d\mathbf{s} = \text{constant}.
\end{equation}

The integral of $\mathbf{u}$ along any closed curve is the \emph{circulation} along the curve. By~\eqref{eq:3-23}, the circulation along any closed curve $C$ remains constant as we follow the curve $C$ with the motion of the fluid elements constituting it.

It is customary to derive the constancy of the circulation along a closed curve first, and then deduce from it the constancy of the flux of the vorticity across a surface. We have followed a reverse procedure to make comparisons with the motions of the magnetic lines of force in hydromagnetics (Chapter~IV, \S\,38\,(b)).

A further important fact is that vortex lines retain their character as vortex lines as we follow the fluid motion. This follows from equation~\eqref{eq:3-13} written somewhat differently. Since
\begin{equation}\label{eq:3-24}
\epsilon_{ijk}\,\epsilon_{klm}\, u_j\, \omega_m = (\delta_{il}\,\delta_{jm} - \delta_{im}\,\delta_{jl})\frac{\partial}{\partial x_j}(u_j\,\omega_i - u_i\,\omega_j) = \frac{\partial}{\partial x_j}(u_j\,\omega_i - u_i\,\omega_j),
\end{equation}
we can write
\begin{equation}\label{eq:3-25}
\frac{\partial\omega_i}{\partial t} + \frac{\partial}{\partial x_j}(u_j\,\omega_i - u_i\,\omega_j) = 0.
\end{equation}

Using the solenoidal character\dag\ of $\mathbf{u}$ and $\boldsymbol{\omega}$, we can rewrite this last equation in the form
\begin{equation}\label{eq:3-26}
\frac{\partial\omega_i}{\partial t} + u_j\frac{\partial\omega_i}{\partial x_j} = \omega_j\frac{\partial u_i}{\partial x_j},
\end{equation}
or
\begin{equation}\label{eq:3-27}
\frac{\partial\boldsymbol{\omega}}{\partial t} + (\mathbf{u}\cdot\operatorname{grad})\boldsymbol{\omega} = \frac{d\boldsymbol{\omega}}{dt} = (\boldsymbol{\omega}\cdot\operatorname{grad})\mathbf{u}.
\end{equation}

Consider now a vortex line. Let $A$ and $B$ be two adjacent points on it. Since the element of arc joining $A$ and $B$ is in the direction of $\boldsymbol{\omega}$,
\begin{equation}\label{eq:3-28}
\overrightarrow{AB} = \delta q\,\boldsymbol{\omega},
\end{equation}
where $\delta q$ is an infinitesimal constant. The velocities at $A$ and $B$ differ by the amount
\begin{equation}\label{eq:3-29}
\mathbf{u}_B - \mathbf{u}_A = (\overrightarrow{AB}\cdot\operatorname{grad})\mathbf{u}_A = \delta q\,(\boldsymbol{\omega}\cdot\operatorname{grad})\mathbf{u}_A;
\end{equation}
or, making use of equation~\eqref{eq:3-27}, we have
\begin{equation}\label{eq:3-30}
\mathbf{u}_B - \mathbf{u}_A = \delta q\,\frac{d\boldsymbol{\omega}}{dt}.
\end{equation}

During a time $\Delta t$, the points $A$ and $B$ will suffer displacements of amounts $\mathbf{u}_A\,\Delta t$ and $\mathbf{u}_B\,\Delta t$, respectively. If the displaced positions of $A$ and $B$ are $A'$ and $B'$, then
\begin{equation}\label{eq:3-31}
\overrightarrow{A'B'} = \overrightarrow{AB} + (\mathbf{u}_B - \mathbf{u}_A)\Delta t
= \delta q\!\left(\boldsymbol{\omega} + \frac{d\boldsymbol{\omega}}{dt}\,\Delta t\right) = \delta q\,\boldsymbol{\omega}_{t+\Delta t}.
\end{equation}

In other words, $\overrightarrow{A'B'}$ is tangential to the vortex line through $A'$ at time $t+\Delta t$. \emph{A vortex line, therefore, consists of the same fluid elements: it moves with the fluid like a material substance.} We may if we like say: \emph{the vortex lines are permanently attached to the fluid.} If a vortex line should disappear in a real fluid, it can only be because of dissipation by viscosity.

\section{The equations of hydrodynamics in a rotating frame of reference}\label{sec:3-21}
Consider a fluid in rotation about some fixed axis with a constant angular velocity $\boldsymbol{\Omega}$. It will be convenient to describe the motions which occur in it as they will appear to an observer at rest in a frame rotating about the same axis and with the same angular velocity. In such a rotating frame of reference, quantities which will be recognized as velocities and accelerations by an observer at rest in it are not the same as the velocities and accelerations as recognized by an observer at rest with respect to the fixed inertial frame $(\xi, \eta, \zeta)$ with respect to which the chosen coordinate system $(x,y,z)$ is rotating with the angular velocity $\Omega$ about the $\zeta$-axis. Also, the $z$- and the $\zeta$-axes will be assumed to coincide. The transformation relating the two coordinate systems is
\begin{align}
x &= +\xi\cos\Omega t + \eta\sin\Omega t, \label{eq:3-32}\\
y &= -\xi\sin\Omega t + \eta\cos\Omega t, \notag\\
z &= +\zeta. \notag
\end{align}

A vector $\mathbf{q}$, with components $q_\xi$, $q_\eta$, and $q_\zeta$ in the inertial frame, will have the following components along the instantaneous directions of the axes of the rotating frame:
\begin{align}
q_x^{(0)} &= +q_\xi\cos\Omega t + q_\eta\sin\Omega t, \label{eq:3-33}\\
q_y^{(0)} &= -q_\xi\sin\Omega t + q_\eta\cos\Omega t, \notag\\
q_z^{(0)} &= +q_\zeta. \notag
\end{align}

We have distinguished the components of the vector obtained by this transformation by a superscript (0) because they are not necessarily the same quantities which will be recognized by the observer, at rest in the rotating frame, as having the meaning of $\mathbf{q}$. The superscript (0) will distinguish the quantities which will have the same \emph{absolute}, as distinct from \emph{relative}, meanings. The reason for this distinction will become clear when we consider the quantities which have the meanings of velocity and acceleration in the rotating frame and their relations to the `absolute' velocity and acceleration.

By differentiating equation~\eqref{eq:3-32} with respect to $t$, we obtain
\begin{align}
\frac{dx}{dt} &= \left(+\frac{d\xi}{dt}\cos\Omega t + \frac{d\eta}{dt}\sin\Omega t\right) - \Omega(\xi\sin\Omega t - \eta\cos\Omega t), \label{eq:3-34}\\
\frac{dy}{dt} &= \left(-\frac{d\xi}{dt}\sin\Omega t + \frac{d\eta}{dt}\cos\Omega t\right) - \Omega(\xi\cos\Omega t + \eta\sin\Omega t). \label{eq:3-35}
\end{align}

Clearly, $dx/dt$ and $dy/dt$ are the quantities which will be recognized by the observer at rest in the rotating frame as the components $u_x$ and $u_y$ of the velocities of the fluid element along the $x$- and the $y$-directions; whereas the quantities grouped together in the first brackets on the right-hand sides of equations~\eqref{eq:3-34} and~\eqref{eq:3-35} are the components along the instantaneous directions of $x$ and $y$ of what is the velocity in the inertial frame; they are the quantities which must be distinguished by the superscript (0). Equations~\eqref{eq:3-34} and~\eqref{eq:3-35} relate $u_x$ and $u_y$ with $u_x^{(0)}$ and $u_y^{(0)}$; we have
\begin{equation}\label{eq:3-36}
u_x = u_x^{(0)} + \Omega y; \quad u_y = u_y^{(0)} - \Omega x; \quad u_z = u_z^{(0)}.
\end{equation}

These equations can be combined into the single vector equation
\begin{equation}\label{eq:3-37}
\mathbf{u} = \mathbf{u}^{(0)} - \boldsymbol{\Omega}\times\mathbf{r}.
\end{equation}

Now differentiating equations~\eqref{eq:3-34} and~\eqref{eq:3-35} once again with respect to $t$, we obtain
\begin{align}
\frac{d^2 x}{dt^2} &= \left(+\frac{d^2\xi}{dt^2}\cos\Omega t + \frac{d^2\eta}{dt^2}\sin\Omega t\right) + 2\Omega\!\left(-\frac{d\xi}{dt}\sin\Omega t + \frac{d\eta}{dt}\cos\Omega t\right) - \Omega^2 x, \label{eq:3-38}\\
\frac{d^2 y}{dt^2} &= \left(-\frac{d^2\xi}{dt^2}\sin\Omega t + \frac{d^2\eta}{dt^2}\cos\Omega t\right) + 2\Omega\!\left(-\frac{d\xi}{dt}\cos\Omega t - \frac{d\eta}{dt}\sin\Omega t\right) - \Omega^2 y. \notag
\end{align}

The equations are equivalent to
\begin{align}
\frac{d^2 x}{dt^2} &= \left(\frac{du_x^{(0)}}{dt}\right)^{(0)} + 2\Omega u_y^{(0)} - \Omega^2 x, \label{eq:3-39}\\
\frac{d^2 y}{dt^2} &= \left(\frac{du_y^{(0)}}{dt}\right)^{(0)} - 2\Omega u_x^{(0)} - \Omega^2 y. \notag
\end{align}

Substituting for $u_x^{(0)}$ and $u_y^{(0)}$ from~\eqref{eq:3-36} in the foregoing equations, we obtain
\begin{align}
\left(\frac{du_x^{(0)}}{dt}\right)^{(0)} &= \frac{du_x}{dt} - 2\Omega u_y - \Omega^2 x, \label{eq:3-40}\\
\left(\frac{du_y^{(0)}}{dt}\right)^{(0)} &= \frac{du_y}{dt} + 2\Omega u_x - \Omega^2 y, \notag\\
\left(\frac{du_z^{(0)}}{dt}\right)^{(0)} &= \frac{du_z}{dt}. \notag
\end{align}

These equations can be combined into the single vector equation
\begin{equation}\label{eq:3-41}
\left(\frac{d\mathbf{u}^{(0)}}{dt}\right)^{(0)} = \frac{d\mathbf{u}}{dt} + 2\boldsymbol{\Omega}\times\mathbf{u} - \tfrac{1}{2}\operatorname{grad}(|\boldsymbol{\Omega}\times\mathbf{r}|^2).
\end{equation}

The term $2\boldsymbol{\Omega}\times\mathbf{u}$ in this equation represents the Coriolis acceleration and the term $-\frac{1}{2}\operatorname{grad}(|\boldsymbol{\Omega}\times\mathbf{r}|^2)$ represents the centrifugal force.

For differentiations with respect to the space coordinates, there are no complicating considerations such as we have just explained for differentiations with respect to time leading to velocities and accelerations. Accordingly, the standard hydrodynamical equation (equation~II~(17)),
\begin{equation}\label{eq:3-42}
\frac{du_i}{dt} = X_i + \frac{1}{\rho}\frac{\partial P_{ij}}{\partial x_j}
\end{equation}
(where $X_i$ is the $i$th component of whatever external forces may be acting on the fluid and $P_{ij}$ is the stress tensor), in the rotating frame of reference, becomes
\begin{equation}\label{eq:3-43}
\frac{du_i}{dt} + u_j\frac{\partial u_i}{\partial x_j} = X_i + \frac{1}{\rho}\frac{\partial P_{ij}}{\partial x_j} + \frac{\partial}{\partial x_i}\!\left(\tfrac{1}{2}|\boldsymbol{\Omega}\times\mathbf{r}|^2\right) + 2\epsilon_{ijk}\, u_j\,\Omega_k.
\end{equation}

The equation of continuity and the equation of heat conduction remain unaffected.

For an incompressible fluid, the equations of motion take the explicit form
\begin{equation}\label{eq:3-44}
\frac{\partial u_i}{\partial t} + u_j\frac{\partial u_i}{\partial x_j} = X_i - \frac{\partial}{\partial x_i}\!\left(\frac{p}{\rho} - \tfrac{1}{2}|\boldsymbol{\Omega}\times\mathbf{r}|^2\right) + \nu\nabla^2 u_i + 2\epsilon_{ijk}\, u_j\,\Omega_k.
\end{equation}

The equation of continuity in the form~\eqref{eq:3-2} continues to hold.

\section{The Taylor--Proudman theorem}\label{sec:3-22}
Consider the equation of motion~\eqref{eq:3-44} for an inviscid fluid in case the external forces are derivable from a potential $V$. We then have
\begin{equation}\label{eq:3-45}
\frac{\partial u_i}{\partial t} + u_j\frac{\partial u_i}{\partial x_j} = -\frac{\partial}{\partial x_i}\!\left(\frac{p}{\rho} - \tfrac{1}{2}|\boldsymbol{\Omega}\times\mathbf{r}|^2 + V\right) + 2\epsilon_{ijk}\, u_j\,\Omega_k.
\end{equation}

Making use of the identity~\eqref{eq:3-9}, we can write
\begin{equation}\label{eq:3-46}
\frac{\partial u_i}{\partial t} - \epsilon_{ijk}\, u_j\, \omega_k = -\frac{\partial\varpi}{\partial x_i} + 2\epsilon_{ijk}\, u_j\,\Omega_k,
\end{equation}
where now
\begin{equation}\label{eq:3-47}
\varpi = p/\rho + \tfrac{1}{2}|\mathbf{u}|^2 + V - \tfrac{1}{2}|\boldsymbol{\Omega}\times\mathbf{r}|^2.
\end{equation}

Alternatively, we can also write
\begin{equation}\label{eq:3-48}
\frac{\partial\mathbf{u}}{\partial t} - \mathbf{u}\times(\boldsymbol{\omega}+2\boldsymbol{\Omega}) = -\operatorname{grad}\varpi,
\end{equation}
and taking the curl of this equation, we have
\begin{equation}\label{eq:3-49}
\frac{\partial\boldsymbol{\omega}}{\partial t} - \operatorname{curl}[\mathbf{u}\times(\boldsymbol{\omega}+2\boldsymbol{\Omega})] = 0.
\end{equation}

Since $\boldsymbol{\Omega}$ is a constant vector, it is clear from a comparison of equations~\eqref{eq:3-13} and~\eqref{eq:3-49} that \emph{in a rotating frame of reference $\boldsymbol{\omega}+2\boldsymbol{\Omega}$ plays the role of\/ $\boldsymbol{\omega}$ in an inertial frame}. We can, therefore, conclude that now
\begin{equation}\label{eq:3-50}
\int_S (\boldsymbol{\omega}+2\boldsymbol{\Omega})\cdot d\mathbf{S} = \text{constant}
\end{equation}
as we follow a surface $S$ (enclosed by a simple closed curve $C$) with the motion in the rotating frame. An equivalent statement is:

\emph{The circulation round a closed contour $C$ bounding a surface $S$ $+$ $2\boldsymbol{\Omega}\times$(the projection of the area of $S$ on a plane perpendicular to $\boldsymbol{\Omega}$)}
\begin{equation}\label{eq:3-51}
= \text{constant}.
\end{equation}

In this form, the theorem appears to have been first stated by Taylor. When the motions are slow and steady, the theorem asserts:

\begin{equation}\label{eq:3-52}
\textit{The projection of the area of\/ $S$ on a plane perpendicular to\/ $\boldsymbol{\Omega}$ is a constant.}
\end{equation}

What this implies for the motions themselves can be inferred more directly from equation~\eqref{eq:3-49}. In a state of steady, slow motions when the squares of the velocity components can be neglected, equation~\eqref{eq:3-49} requires that
\begin{equation}\label{eq:3-53}
\operatorname{curl}(\mathbf{u}\times\boldsymbol{\Omega}) = 0,
\end{equation}
or (cf.\ equation~\eqref{eq:3-24})
\begin{equation}\label{eq:3-54}
\frac{\partial}{\partial x_j}(\omega_i\,\Omega_j - u_j\,\Omega_i) = 0.
\end{equation}

Since $\boldsymbol{\Omega}$ is a constant vector and $\mathbf{u}$ is solenoidal, equation~\eqref{eq:3-54} reduces to
\begin{equation}\label{eq:3-55}
\Omega_j\frac{\partial u_i}{\partial x_j} = 0.
\end{equation}

\emph{The motions cannot, therefore, vary in the direction of\/ $\boldsymbol{\Omega}$. In other words, all steady slow motions in a rotating inviscid fluid are necessarily two-dimensional.} This is the form in which the Taylor--Proudman theorem is generally stated; equation~\eqref{eq:3-50} is, however, a more complete statement.

What the condition~\eqref{eq:3-55} implies is this: the fact that the motions transverse to $\boldsymbol{\Omega}$ cannot vary in the direction of $\boldsymbol{\Omega}$ means that any two fluid elements which are initially on a line parallel to $\boldsymbol{\Omega}$ will always remain on that line; and the fact that motions in the direction of $\boldsymbol{\Omega}$ cannot also vary along this direction means that two fluid elements which are initially a certain distance apart (in the direction of $\boldsymbol{\Omega}$) will always remain at this distance apart. This predicted behaviour of slow, steady motions in a rotating fluid was strikingly demonstrated by Taylor as follows.

Water in a tank was first made to rotate steadily as a solid body. A small motion was then communicated to it and a few drops of coloured liquid (or ink) were inserted. The slow movement of the fluid then drew this coloured portion of the fluid into sheets; and these sheets

\figurePlaceholder{20}{Illustrating the Taylor--Proudman theorem in a rotating cylinder of water. In the experiment, ink was poured into the rotating cylinder and stirred $1\frac{1}{2}$~min later by dipping a stationary rod for 15~sec. Note the thinness of many of the filaments into which the ink is drawn in accordance with the Taylor--Proudman theorem.}

\noindent always remained parallel to the axis of rotation. Taylor noted: `The accuracy with which they remained parallel to the axis of rotation is quite extraordinary.' These experiments of Taylor have been repeated by Fultz; one of his photographs illustrating this phenomenon is reproduced in Fig.~20.

\section{The propagation of waves in a rotating fluid}\label{sec:3-23}
We shall conclude this excursion into general hydrodynamical theorems by showing that the equations of motion in a rotating fluid allow periodic solutions representing the propagation of waves.

In the absence of external forces, or in the presence of forces which are derivable from a potential, the equations of motion governing a viscous incompressible fluid take the form (cf.\ equation~\eqref{eq:3-44})
\begin{equation}\label{eq:3-56}
\frac{\partial u_i}{\partial t} + u_j\frac{\partial u_i}{\partial x_j} = -\frac{\partial\varpi}{\partial x_i} + \nu\nabla^2 u_i + 2\epsilon_{ijk}\, u_j\,\Omega_k.
\end{equation}

We seek solutions of this equation with a space-time dependence given by
\begin{equation}\label{eq:3-57}
e^{i(\mathbf{k}\cdot\mathbf{r}+pt)}.
\end{equation}

For solutions having this space-time dependence,
\begin{equation}\label{eq:3-58}
\frac{\partial}{\partial t} = ip, \quad \frac{\partial}{\partial x_j} = ik_j, \quad \text{and} \quad \nabla^2 = -k^2;
\end{equation}
and equation~\eqref{eq:3-56} gives
\begin{equation}\label{eq:3-59}
(ip + \nu k^2)u_i + ik_j\, u_j\, u_i = -ik_i\,\varpi + 2\epsilon_{ijk}\, u_j\,\Omega_k;
\end{equation}
while the equation of continuity gives
\begin{equation}\label{eq:3-60}
k_j\, u_j = 0.
\end{equation}

This last equation implies that \emph{the waves are transverse}.

In view of~\eqref{eq:3-60}, equation~\eqref{eq:3-59} becomes
\begin{equation}\label{eq:3-61}
n u_i + k_j\,\varpi + 2i\epsilon_{ijk}\, u_j\,\Omega_k = 0,
\end{equation}
where we have written
\begin{equation}\label{eq:3-62}
n = p - i\nu k^2.
\end{equation}

Multiplying equation~\eqref{eq:3-61} scalarly by $u_i$, $k_i$, and $\Omega_i$, we obtain, respectively,
\begin{equation}\label{eq:3-63}
\begin{aligned}
n u_i^2 &= 0, \\
k^2\,\varpi + 2i\epsilon_{ijk}\, k_i\, u_j\,\Omega_k &= 0, \\
n u_i\,\Omega_i + \varpi k_i\,\Omega_i &= 0.
\end{aligned}
\end{equation}

It is now convenient to choose the orientation of the coordinate system such that
\begin{equation}\label{eq:3-64}
\boldsymbol{\Omega} = (\Omega_x, 0, \Omega_z) \quad \text{and} \quad \mathbf{k} = (0,0,k);
\end{equation}
this clearly implies no loss of generality. With this choice of the coordinate system, equation~\eqref{eq:3-60} gives
\begin{equation}\label{eq:3-65}
k u_z = 0 \quad \text{or} \quad u_z = 0.
\end{equation}

Equations~\eqref{eq:3-63} now give
\begin{align}
n(u_x^2 + u_y^2) &= 0, \label{eq:3-66}\\
2\Omega_z\, u_y &= +k\varpi, \label{eq:3-67}\\
n\Omega_x\, u_x &= -k\varpi\Omega_z. \label{eq:3-68}
\end{align}

These equations enable us to express the amplitudes $u_x$ and $\varpi$ in terms of $u_z$; we find
\begin{equation}\label{eq:3-69}
u_y = +i\frac{n}{2\Omega_z}u_x; \quad \varpi = -\frac{n\Omega_x}{k\Omega_z}u_x; \quad u_z = 0.
\end{equation}

Substituting the first of these relations in equation~\eqref{eq:3-66}, we obtain
\begin{equation}\label{eq:3-70}
n\!\left(1 - \frac{n^2}{4\Omega_z^2}\right)\!u_x^2 = 0.
\end{equation}

From this equation it follows that $n$ is real; and ignoring the root $n = 0$, we have
\begin{equation}\label{eq:3-71}
n = \pm 2\Omega_z = \pm 2\Omega\cos\vartheta,
\end{equation}
where $\vartheta$ is the inclination of the direction of wave-propagation to the direction of $\boldsymbol{\Omega}$. The gyration frequency $p$ is given by (cf.\ equation~\eqref{eq:3-62})
\begin{equation}\label{eq:3-72}
p = n + i\nu k^2 = \pm 2\Omega\cos\vartheta + i\nu k^2.
\end{equation}

The waves are, therefore, damped; and the damping constant is $\nu k^2$.

With $n$ given by equation~\eqref{eq:3-72}, the expressions for the amplitudes become
\begin{equation}\label{eq:3-73}
u_x = \pm i u_y, \quad \varpi = \mp\frac{2\Omega}{k}u_x\sin\vartheta, \quad u_z = 0.
\end{equation}

\emph{The waves are, therefore, transverse, circularly polarized, and damped.}

In the limit of zero viscosity, the phase velocity is given by
\begin{equation}\label{eq:3-74}
V = \frac{p}{k} = \pm\frac{2\Omega}{k}\cos\vartheta.
\end{equation}

This is the \emph{phase velocity}. The \emph{group velocity} is given by
\begin{equation}\label{eq:3-75}
\frac{\partial p}{\partial k_i} = \pm 2\frac{\partial}{\partial k_i}\!\left(\frac{\Omega_j k_j}{k}\right) = \pm 2\!\left(\frac{\Omega_i}{k} - \frac{\Omega_j k_j}{k^3}k_i\right),
\end{equation}
or
\begin{equation}\label{eq:3-76}
\frac{\partial p}{\partial \mathbf{k}} = \pm\frac{2}{k^3}\mathbf{k}\times(\boldsymbol{\Omega}\times\mathbf{k}).
\end{equation}

Finally, it should be noted that the foregoing solutions representing the propagation of waves are not limited to infinitesimal amplitudes.

\section{The problem of thermal instability in a rotating fluid: general considerations}\label{sec:3-24}
From the theorems proved in \S\S\,22 and~23 it is apparent that rotation should have a profound effect on the onset of thermal instability. For convection implies that motions which occur have, necessarily, a three-dimensional character; this the Taylor--Proudman theorem expressly forbids for an inviscid fluid so long as the non-linear terms in the equations of motion are neglected and the motions are steady. The first of these two conditions is certainly met in a linear stability theory; and the second obtains for steady convection. Therefore, in contrast to non-rotating fluids, an inviscid fluid in rotation should be expected to be thermally stable for \emph{all} adverse temperature gradients. Indeed, only in the presence of viscosity can thermal instability arise; for only then can the Taylor--Proudman theorem be violated.

There is another factor to remember. While the Taylor--Proudman theorem forbids any variation of the velocity in the direction of $\boldsymbol{\Omega}$, it does not forbid oscillatory motions. As we have seen in \S\,23, the propagation of transverse waves is possible under the same circumstances. For this reason, we may anticipate that overstability as a means of developing instability may, under certain conditions, play a decisive role. We shall see that this is indeed the case.

\section{The perturbation equations}\label{sec:3-25}
Consider, then, an infinite horizontal layer of fluid which is kept rotating at a constant rate. Let $\boldsymbol{\Omega}$ denote the angular velocity of rotation. As we have shown explicitly in Chapter~II, \S\,9 for problems of the type we are considering, it will suffice to work with the relevant equations of motion and heat conduction in the Boussinesq approximation. The only additional factors which we have to allow for now are the effects of Coriolis acceleration and centrifugal force in the equations of motion. Thus, in view of equation~\eqref{eq:3-41}, we replace equation~II~(43) by
\begin{equation}\label{eq:3-77}
\frac{\partial u_i}{\partial t} + u_j\frac{\partial u_i}{\partial x_j} = -\frac{\partial}{\partial x_i}\!\left(\frac{\delta p}{\rho_0} - \tfrac{1}{2}|\boldsymbol{\Omega}\times\mathbf{r}|^2\right) + \left(1+\frac{\delta\rho}{\rho_0}\right)X_i + \nu\nabla^2 u_i + 2\epsilon_{ijk}\, u_j\,\Omega_k.
\end{equation}

The remaining equations (equations~II~(41), (44), and (46)) are unaffected.

We envisage now an initial state in which a steady adverse temperature gradient $\beta$ is maintained and there are no motions. The characterization of this initial state differs in no respect from that described in Chapter~II, \S\,9; and the equations governing small perturbations can be obtained in exactly the same way. Thus, equation~II~(55) is replaced by
\begin{equation}\label{eq:3-78}
\frac{\partial u_i}{\partial t} = -\frac{\partial}{\partial x_i}\!\left(\frac{\delta p}{\rho_0}\right) + g\alpha\Theta\,\lambda_i + \nu\nabla^2 u_i + 2\epsilon_{ijk}\, u_j\,\Omega_k,
\end{equation}
where $\boldsymbol{\lambda} = (0,0,1)$ is a unit vector in the direction of the vertical. We also have (equations~II~(56) and~(57))
\begin{equation}\label{eq:3-79}
\frac{\partial\Theta}{\partial t} = \beta\lambda_j u_j + \kappa\nabla^2\Theta
\end{equation}
and
\begin{equation}\label{eq:3-80}
\frac{\partial u_j}{\partial x_j} = 0.
\end{equation}

As in Chapter~II, \S\,9, we can eliminate the term in $\delta p/\rho_0$ in equation~\eqref{eq:3-78} by applying the operator $\epsilon_{ijk}\,\partial/\partial x_j$ to the $k$-component of the equation; since (cf.\ equation~\eqref{eq:3-24})
\begin{equation}\label{eq:3-81}
\epsilon_{ijk}\frac{\partial}{\partial x_j}\epsilon_{klm}\, u_l\,\Omega_m = \frac{\partial}{\partial x_j}(u_i\,\Omega_j - u_j\,\Omega_i) = \Omega_j\frac{\partial u_i}{\partial x_j},
\end{equation}
the result is (cf.\ equation~II~(67))
\begin{equation}\label{eq:3-82}
\frac{\partial\omega_i}{\partial t} = g\alpha\epsilon_{ijk}\frac{\partial\Theta}{\partial x_j}\lambda_k + \nu\nabla^2\omega_i + 2\Omega_j\frac{\partial u_i}{\partial x_j}.
\end{equation}

Taking the curl of this equation once again, we obtain (cf.\ equation~II~(72))
\begin{equation}\label{eq:3-83}
\frac{\partial}{\partial t}(\nabla^2 u_i) = g\alpha\!\left(\lambda_i\nabla^2\Theta - \lambda_j\frac{\partial^2\Theta}{\partial x_i\,\partial x_j}\right) + \nu\nabla^4 u_i - 2\Omega_j\frac{\partial\omega_i}{\partial x_j}.
\end{equation}

Now multiplying equations~\eqref{eq:3-82} and~\eqref{eq:3-83} by $\lambda_i$, we get
\begin{equation}\label{eq:3-84}
\frac{\partial\zeta}{\partial t} = \nu\nabla^2\zeta + 2\Omega_j\frac{\partial w}{\partial x_j},
\end{equation}
and
\begin{equation}\label{eq:3-85}
\frac{\partial}{\partial t}\nabla^2 w = g\alpha\!\left(\frac{\partial^2\Theta}{\partial x^2}+\frac{\partial^2\Theta}{\partial y^2}\right) + \nu\nabla^4 w - 2\Omega_j\frac{\partial\zeta}{\partial x_j},
\end{equation}
where $w$ and $\zeta$ are the $z$-components of the velocity and the vorticity, respectively.

Most of our discussions relating to equations~\eqref{eq:3-84} and~\eqref{eq:3-85} will be restricted to the case when $\boldsymbol{\Omega}$ and $\mathbf{g}$ are parallel; the case when they are not will be briefly considered in \S\,34.

When the axis of rotation coincides with the vertical, the relevant equations are
\begin{align}
\frac{\partial\Theta}{\partial t} &= \beta w + \kappa\nabla^2\Theta, \label{eq:3-86}\\[4pt]
\frac{\partial\zeta}{\partial t} &= 2\Omega\frac{\partial w}{\partial z} + \nu\nabla^2\zeta, \label{eq:3-87}
\end{align}
and
\begin{equation}\label{eq:3-88}
\frac{\partial}{\partial t}\nabla^2 w = g\alpha\!\left(\frac{\partial^2\Theta}{\partial x^2}+\frac{\partial^2\Theta}{\partial y^2}\right) + \nu\nabla^4 w - 2\Omega\frac{\partial\zeta}{\partial z},
\end{equation}
and we have to seek solutions of these equations which satisfy the boundary conditions described in Chapter~II, \S\,9\,(a).

\subsection*{(a) The analysis into normal modes}

Following the procedure in Chapter~II, \S\,10, we analyse the perturbations $w$, $\Theta$, and $\zeta$ into two-dimensional periodic waves; and considering disturbances characterized by a particular wave number $k$, we suppose that $w$, $\Theta$, and $\zeta$ have the forms assumed in equation~II~(90). Equations~\eqref{eq:3-86}--\eqref{eq:3-88} then become
\begin{align}
p\,\Theta &= \beta W + \kappa\!\left(\frac{d^2}{dz^2}-k^2\right)\!\Theta, \label{eq:3-89}\\[4pt]
p\,Z &= 2\Omega\frac{dW}{dz} + \nu\!\left(\frac{d^2}{dz^2}-k^2\right)\!Z, \label{eq:3-90}\\[4pt]
p\!\left(\frac{d^2}{dz^2}-k^2\right)\!W &= -g\alpha k^2\Theta + \nu\!\left(\frac{d^2}{dz^2}-k^2\right)^{\!2}\!W - 2\Omega\frac{dZ}{dz}. \label{eq:3-91}
\end{align}

Measuring lengths in the unit $d$ and letting
\begin{equation}\label{eq:3-92}
a = kd, \quad \sigma = pd^2/\nu, \quad \text{and} \quad \mathscr{P} = \nu/\kappa,
\end{equation}
we can reduce equations~\eqref{eq:3-89}--\eqref{eq:3-91} to the forms
\begin{align}
(D^2 - a^2 - \mathscr{P}\sigma)\,\Theta &= -\left(\frac{\beta d^2}{\kappa}\right)\!W, \label{eq:3-93}\\[4pt]
(D^2 - a^2 - \sigma)\,Z &= -\left(\frac{2\Omega}{\nu}\,d\right)\!DW, \label{eq:3-94}\\[4pt]
(D^2 - a^2)(D^2 - a^2 - \sigma)\,W - \left(\frac{2\Omega}{\nu}\,d^3\right)\!DZ &= \left(\frac{g\alpha d^2}{\nu}\right)\!a^2\Theta, \label{eq:3-95}
\end{align}
where $D = d/dz$ ($z$ being measured in the new unit). Solutions of equations~\eqref{eq:3-93}--\eqref{eq:3-95} must be sought which satisfy the boundary conditions given in equations~II~(95) and~(96).

\section{The case when instability sets in as stationary convection. A variational principle}\label{sec:3-26}

We shall find that in contrast to the simple B\'enard problem, the principle of the exchange of stabilities is not valid, generally, when rotation is present. However, it is not possible to state in simple analytical terms the necessary and sufficient conditions for its validity; it will appear that they are best stated in terms of the solution for the case of instability as stationary convection. For this reason, we shall first establish the conditions for this latter type of instability to occur before deciding whether instability will in fact occur this way.

When instability sets in as ordinary convection, the marginal state will be characterized by $\sigma = 0$, and the basic equations reduce to (cf.\ equations~\eqref{eq:3-93}--\eqref{eq:3-95})
\begin{align}
(D^2 - a^2)\,\Theta &= -\left(\frac{\beta d^2}{\kappa}\right)\!W, \label{eq:3-96}\\[4pt]
(D^2 - a^2)\,Z &= -\left(\frac{2\Omega}{\nu}\,d\right)\!DW, \label{eq:3-97}
\end{align}
and
\begin{equation}\label{eq:3-98}
(D^2 - a^2)^2 W - \left(\frac{2\Omega}{\nu}\,d^3\right)\!DZ = \left(\frac{g\alpha d^2}{\nu}\right)\!a^2\Theta.
\end{equation}

We can eliminate $Z$ and $\Theta$ from the last equation by operating $(D^2 - a^2)$ on it. We find:
\begin{equation}\label{eq:3-99}
(D^2 - a^2)^3 W + \mathscr{T}\,D^2 W = -Ra^2 W,
\end{equation}
where $R$ is the Rayleigh number as usually defined, and
\begin{equation}\label{eq:3-100}
\mathscr{T} = \frac{4\Omega^2}{\nu^2}\,d^4
\end{equation}
is the Taylor number.

The boundary conditions with respect to which we must seek solutions of the foregoing equations are (cf.\ equations~II~(95) and~(96))
\begin{equation}\label{eq:3-101}
W = 0 \quad \text{and} \quad \Theta = 0 \quad \text{for } z = 0 \text{ and } 1,
\end{equation}
\begin{equation}\label{eq:3-102}
\text{either}\quad Z = 0 \quad \text{and} \quad DW = 0 \quad \text{(on a rigid surface)},
\end{equation}
\begin{equation*}
\text{or}\quad DZ = 0 \quad \text{and} \quad D^2 W = 0 \quad \text{(on a free surface)}.
\end{equation*}

In view of equation~\eqref{eq:3-98} we can replace the conditions~\eqref{eq:3-101} by
\begin{equation}\label{eq:3-103}
W = 0 \quad \text{and} \quad (D^2 - a^2)^2 W - \left(\frac{2\Omega}{\nu}\,d^3\right)\!DZ = 0 \quad \text{for } z = 0 \text{ and } 1.
\end{equation}

Since the boundary conditions~\eqref{eq:3-103} involve $Z$, it is clear that we cannot treat equation~\eqref{eq:3-99} independently of equation~\eqref{eq:3-97} for $Z$; the order of the system we are effectively dealing with is, therefore, eight not six.

Equations~\eqref{eq:3-97} and~\eqref{eq:3-99}, together with the boundary conditions~\eqref{eq:3-102} and~\eqref{eq:3-103}, clearly constitute a characteristic value problem for $R$ for given $a^2$ and $\mathscr{T}$; and the problem of determining the critical Rayleigh number for the onset of instability as stationary convection reduces to the following.

We must first determine the lowest characteristic value of $R$, through the solution of equations~\eqref{eq:3-97}, \eqref{eq:3-99}, \eqref{eq:3-102}, and~\eqref{eq:3-103}, as a function of $a^2$ for a given $\mathscr{T}$ and then find the minimum of this function; and the minimum so determined is the critical Rayleigh number for the onset of stationary convection for the given~$\mathscr{T}$.

\subsection*{(a) A variational principle}

As in the case of the simple B\'enard problem, we can formulate the present characteristic value problem also in terms of a variational principle.

Letting
\begin{equation}\label{eq:3-104}
F = (D^2 - a^2)^2 W - \left(\frac{2\Omega}{\nu}\,d^3\right)\!DZ,
\end{equation}
we can rewrite the differential equations governing $W$ and $Z$ in the forms
\begin{equation}\label{eq:3-105}
(D^2 - a^2)F = -Ra^2 W
\end{equation}
and
\begin{equation}\label{eq:3-106}
(D^2 - a^2)Z = -\left(\frac{2\Omega}{\nu}\,d\right)\!DW.
\end{equation}

The boundary conditions~\eqref{eq:3-103} require that
\begin{equation}\label{eq:3-107}
W = F = 0 \quad \text{for } z = 0 \text{ and } 1.
\end{equation}

Now multiply equation~\eqref{eq:3-105} by $F$ and integrate over the range of $z$. After an integration by parts, the left-hand side gives
\begin{equation}\label{eq:3-108}
\int_0^1 F(D^2 - a^2)F\,dz = -\int_0^1 [(DF)^2 + a^2 F^2]\,dz,
\end{equation}
the integrated part vanishing on account of the boundary condition on $F$; and the right-hand side of the equation requires us to consider
\begin{equation}\label{eq:3-109}
\int_0^1 W F\,dz = \int_0^1 W(D^2 - a^2)^2 W\,dz - \frac{2\Omega}{\nu}\,d^3 \int_0^1 W\,DZ\,dz.
\end{equation}

After two integrations by parts, the first of the two integrals on the right-hand side of~\eqref{eq:3-109} becomes (cf.\ equation~II~(120))
\begin{equation}\label{eq:3-110}
\int_0^1 W(D^2 - a^2)^2 W\,dz = \int_0^1 [(D^2 - a^2)W]^2\,dz,
\end{equation}
while the second, after an integration by parts, gives (cf.\ equations~\eqref{eq:3-103} and~\eqref{eq:3-106})
\begin{equation}\label{eq:3-111}
-\frac{2\Omega}{\nu}\,d^3\int_0^1 W\,DZ\,dz = \frac{2\Omega}{\nu}\,d^3\int_0^1 Z\,DW\,dz = -d^2\int_0^1 Z(D^2-a^2)Z\,dz.
\end{equation}

Remembering that on a bounding surface either $Z$ or $DZ$ vanishes, we obtain after a further integration by parts
\begin{equation}\label{eq:3-112}
-\frac{2\Omega}{\nu}\,d^3\int_0^1 W\,DZ\,dz = d^2 \int_0^1 [(DZ)^2 + a^2 Z^2]\,dz.
\end{equation}

Combining equations~\eqref{eq:3-109}, \eqref{eq:3-110}, and~\eqref{eq:3-112}, we have
\begin{equation}\label{eq:3-113}
\int_0^1 W F\,dz = \int_0^1 [(D^2 - a^2)W]^2\,dz + d^2 \int_0^1 [(DZ)^2 + a^2 Z^2]\,dz.
\end{equation}

The result of multiplying equation~\eqref{eq:3-105} by $F$ and integrating over $z$ is
\begin{equation}\label{eq:3-114}
R = \frac{\displaystyle\int_0^1 [(DF)^2 + a^2 F^2]\,dz}{a^2\displaystyle\int_0^1\bigl\{[(D^2 - a^2)W]^2 + d^2[(DZ)^2 + a^2 Z^2]\bigr\}\,dz} = \frac{I_1}{a^2 I_2} \quad \text{(say)}.
\end{equation}

This formula expresses $R$ as a ratio of two positive definite integrals. Consider now the effect on $R$ of variations $\delta W$ and $\delta Z$ in $W$ and $Z$ compatible with the boundary conditions on $W$ and $Z$. To the first order, we have
\begin{equation}\label{eq:3-115}
\delta R = \frac{1}{a^2 I_2}\!\left(\delta I_1 - \frac{I_1}{I_2}\,\delta I_2\right) = \frac{1}{a^2 I_2}(\delta I_1 - Ra^2\,\delta I_2).
\end{equation}

$\delta I_1$ and $\delta I_2$ are the corresponding variations in $I_1$ and $I_2$:
\begin{align}
\delta I_1 &= 2\int_0^1 [(DF)(D\,\delta F) + a^2 F\,\delta F]\,dz, \notag\\[6pt]
\delta I_2 &= 2\int_0^1 \bigl\{(D^2 - a^2)W [(D^2 - a^2)\,\delta W] + \notag\\
&\qquad\qquad + d^2[(DZ)(D\,\delta Z) + a^2 Z\,\delta Z]\bigr\}\,dz. \label{eq:3-116}
\end{align}

Making use of the boundary conditions which $F$, $\delta F$, $W$, $\delta W$, $Z$, and $\delta Z$ satisfy, we can reduce the expressions for $\delta I_1$ and $\delta I_2$ by one or more integrations by parts. Thus
\begin{equation}\label{eq:3-117}
\delta I_1 = -2\int_0^1 \delta F(D^2 - a^2)F\,dz,
\end{equation}
and
\begin{align}
\delta I_2 &= 2\int_0^1 W(D^2 - a^2)^2\,\delta W - 2d^2\int_0^1 \delta Z(D^2 - a^2)Z\,dz \notag\\
&= 2\int_0^1 W[(D^2 - a^2)^2\,\delta W + \frac{4\Omega^2}{\nu^2}d^4\int_0^1 \delta Z\,DW\,dz \notag\\
&= 2\int_0^1 W\!\left[(D^2 - a^2)^2\,\delta W - \left(\frac{2\Omega}{\nu}\,d^3\right)\!D\,\delta Z\right]\!dz \label{eq:3-118}\\
&= 2\int_0^1 W\,\delta F\,dz. \notag
\end{align}

Now combining equations~\eqref{eq:3-115}, \eqref{eq:3-117}, and~\eqref{eq:3-118}, we have
\begin{equation}\label{eq:3-119}
\delta R = -\frac{2}{a^2 I_2}\int_0^1 \delta F[(D^2 - a^2)F + Ra^2 W]\,dz,
\end{equation}
where it may be recalled that equation~\eqref{eq:3-104} defining $F$ and equation~\eqref{eq:3-106} relating $Z$ and $W$ have been explicitly used in the reductions.

From equation~\eqref{eq:3-119} it follows that
\begin{equation}\label{eq:3-120}
\delta R = 0 \quad \text{if}\quad (D^2 - a^2)F = -Ra^2 W;
\end{equation}
\emph{and conversely if $\delta R = 0$ for any arbitrary variation,}
\[
\delta F = (D^2 - a^2)^2\,\delta W - (2\Omega d^3/\nu)\,D\,\delta Z,
\]
\emph{compatible with the boundary conditions of the problem, then equation~\eqref{eq:3-120} must hold and the function $W$ in terms of which $R$ was initially calculated must have been a solution of the characteristic value problem.}

The foregoing arguments establish the stationary property of the characteristic values when they are regarded as given by equation~\eqref{eq:3-114}. And it can be shown, by a procedure exactly analogous to that used in connexion with the first variational principle for the simple B\'enard problem (Chapter~II, \S\,13\,(a)), that the lowest characteristic value of $R$ is, indeed, a true minimum of the functional~\eqref{eq:3-114}. We are thus enabled to formulate the following variational procedure for solving equations~\eqref{eq:3-97} and~\eqref{eq:3-99} (for any assigned $a^2$ and $\mathscr{T}$) and satisfying the boundary conditions of the problem.

Assume for $F$ an expansion involving one or more parameters $A_p$ which vanishes for $z = 0$ and~1. With the chosen form of $F$, determine $W$ and $Z$ as solutions of the equations
\begin{equation}\label{eq:3-121}
(D^2 - a^2)^2 W - \left(\frac{2\Omega}{\nu}\,d^3\right)\!DZ = F
\end{equation}
and
\begin{equation}\label{eq:3-122}
(D^2 - a^2)Z = -\left(\frac{2\Omega}{\nu}\,d\right)\!DW,
\end{equation}
which satisfy a total of six boundary conditions on $W$ and $Z$ at $z = 0$ and~1. Since equations~\eqref{eq:3-121} and~\eqref{eq:3-122} are together of order six, there will be just enough constants of integration in the general solution to satisfy all the boundary conditions. Having solved for $W$ and $Z$ in this manner, evaluate $R$ according to the formula~\eqref{eq:3-114} and minimize it with respect to the parameters $A_p$. In this way, we shall obtain the `best' value of $R$ for the chosen form of $F$. We shall see in the following section that even with the simplest trial function for $F$, we can reach quite high precision in the deduced values of $R$.

Finally, it may be noted that according to equations~\eqref{eq:3-121} and~\eqref{eq:3-122}, the equation for $W$ that must be solved when using the variational method is
\begin{equation}\label{eq:3-123}
(D^2 - a^2)^3 W + \mathscr{T}\,D^2 W = (D^2 - a^2)F,
\end{equation}
where $\mathscr{T}$ is the Taylor number.

\section{Solutions for the case when instability sets in as stationary convection}\label{sec:3-27}

We shall obtain the solution for the following three cases: when both the bounding surfaces are free; when both the bounding surfaces are rigid; and when one bounding surface is rigid and the other is free.

\subsection*{(a) The solution for two free boundaries}

In this case the boundary conditions~\eqref{eq:3-102} and~\eqref{eq:3-103} require
\begin{equation}\label{eq:3-124}
W = D^2 W = D^4 W = 0 \quad \text{and} \quad DZ = 0 \quad \text{for } z = 0 \text{ and } 1.
\end{equation}

From the equation satisfied by $W$ (namely~\eqref{eq:3-99}), it follows that $D^6 W = 0$ for $z = 0$ and~1. By differentiating equation~\eqref{eq:3-99} even numbers of times, we can successively conclude that all the even derivatives of $W$ must, therefore, be
\begin{equation}\label{eq:3-125}
W = A\sin n\pi z,
\end{equation}
where $A$ is a constant and $n$ is an integer. The corresponding solution for $Z$ is
\begin{equation}\label{eq:3-126}
Z = A\!\left(\frac{2\Omega}{\nu}\,d\right)\!\frac{n\pi}{n^2\pi^2 + a^2}\cos n\pi z.
\end{equation}

Substitution of the solution~\eqref{eq:3-125} in equation~\eqref{eq:3-99} leads to the characteristic equation
\begin{equation}\label{eq:3-127}
R = \frac{1}{a^2}[(n^2\pi^2 + a^2)^3 + n^2\pi^2\,\mathscr{T}].
\end{equation}

For a given $a^2$, the lowest characteristic value for $R$ occurs when $n = 1$; then,
\begin{equation}\label{eq:3-128}
R = \frac{1}{a^2}[(\pi^2 + a^2)^3 + \pi^2\,\mathscr{T}];
\end{equation}
and this equation has to be interpreted in the same way as equation~II~(193) was interpreted in Chapter~II, \S\,15\,(a).

Letting
\begin{equation}\label{eq:3-129}
a^2 = \pi^2 x,
\end{equation}
we can rewrite equation~\eqref{eq:3-128} in the form
\begin{equation}\label{eq:3-130}
R = \pi^4\!\left[(1+x)^3 + \frac{\mathscr{T}}{\pi^4}\right]\!\frac{1}{x}.
\end{equation}

\begin{table}[ht]
\centering
\caption{Critical Rayleigh numbers and wave numbers of the unstable modes at marginal stability for the onset of stationary convection when both bounding surfaces are free}
\label{tab:VII}
\begin{tabular}{ccccccc}
\hline
$\mathscr{T}$ & $a_c$ & $R_c$ & $\mathscr{T}$ & $a_c$ & $R_c$ \\
\hline
0 & $2{\cdot}233$ & $6{\cdot}575\times 10^2$ & $3\times 10^4$ & $10{\cdot}45$ & $4{\cdot}237\times 10^4$ \\
$10$ & $2{\cdot}270$ & $6{\cdot}771\times 10^2$ & $10^5$ & $12{\cdot}85$ & $9{\cdot}232\times 10^4$ \\
$10^2$ & $2{\cdot}594$ & $8{\cdot}263\times 10^2$ & $10^6$ & $19{\cdot}02$ & $4{\cdot}147\times 10^5$ \\
$5\times 10^2$ & $3{\cdot}228$ & $1{\cdot}473\times 10^3$ & $10^7$ & $28{\cdot}02$ & $1{\cdot}897\times 10^6$ \\
$10^3$ & $3{\cdot}710$ & & $10^8$ & $41{\cdot}20$ & $8{\cdot}536\times 10^6$ \\
$10^4$ & $4{\cdot}211$ & $2{\cdot}009\times 10^4$ & $10^9$ & $60{\cdot}53$ & $4{\cdot}017\times 10^7$ \\
$5\times 10^3$ & $5{\cdot}021$ & $5{\cdot}999\times 10^3$ & $10^{10}$ & $88{\cdot}87$ & $1{\cdot}879\times 10^8$ \\
$5\times 10^3$ & $5{\cdot}665$ & $3{\cdot}377\times 10^4$ & $10^{11}$ & $130{\cdot}46$ & $8{\cdot}702\times 10^8$ \\
$3\times 10^4$ & $6{\cdot}054$ & $1{\cdot}013\times 10^4$ & $10^{12}$ & $191{\cdot}31$ & $4{\cdot}037\times 10^9$ \\
$10^4$ & $8{\cdot}245$ & & & & \\
\hline
\end{tabular}
\end{table}

As a function of $x$, $R$ given by equation~\eqref{eq:3-130} attains its minimum when
\begin{equation}\label{eq:3-131}
2x^3 + 3x^2 = 1 + \mathscr{T}/\pi^4.
\end{equation}

With $x$ determined as the root of this cubic equation, equation~\eqref{eq:3-130} will give the required critical Rayleigh number $R_c$. Values of $R_c$ determined in this fashion for various values of $\mathscr{T}$ are given in Table~VII; the wave numbers characterizing the marginal states are also given. The results are further illustrated in Figs.~21 and~22. The inhibiting effect of rotation on the onset of instability is apparent from these results.

For $\mathscr{T}/\pi^4$ sufficiently large, the required root of equation~\eqref{eq:3-131} tends to
\begin{equation}\label{eq:3-132}
x_{\min} \to \left(\frac{\mathscr{T}}{2\pi^4}\right)^{\!1/3} \quad (\mathscr{T}\to\infty).
\end{equation}

The corresponding asymptotic behaviours of $R_c$ and $a_{\min}$ are:
\begin{equation}\label{eq:3-133}
R_c \to 3\pi^4\!\left(\frac{\mathscr{T}}{2\pi^4}\right)^{\!2/3} = 8{\cdot}6956\,\mathscr{T}^{2/3} \quad (\mathscr{T}\to\infty).
\end{equation}
and
\[
a_{\min} \to (\tfrac{1}{2}\pi^2\,\mathscr{T})^{1/6} = 1{\cdot}3048\,\mathscr{T}^{1/6}.
\]

\figurePlaceholder{21}{The $(R_c,\mathscr{T})$-relations for the three cases (i) both bounding surfaces rigid, (ii) one bounding surface rigid and the other free, and (iii) both bounding surfaces free; the curves labelled $a$, $b$, and $c$ are the relations for the onset of ordinary cellular convection for the three cases, respectively. The curves labelled $a'A$, $b'B$, and $c'C'$ are the corresponding relations for the onset of overstability for $\mathscr{P} = 0{\cdot}025$. At $c$ (respectively $c'$) we have a change from one type of instability to another as $\mathscr{T}$ increases.}

\figurePlaceholder{22}{The $(a_c,\mathscr{T})$-relations for the three cases (i) both bounding surfaces rigid, (ii) one bounding surface rigid and the other free, and (iii) both bounding surfaces free; the curves labelled $a$, $b$, and $c$ are the relations for the onset of ordinary cellular convection for the three cases, respectively. The curves labelled $a'A$, $B$, and $C'$ are the corresponding relations for the onset of overstability for $\mathscr{P} = 0{\cdot}025$.}

Substituting for $R$ and $\mathscr{T}$ in accordance with their definitions, we find that the formula which determines the critical temperature gradient for the onset of stationary convection, when $\mathscr{T}\to\infty$, is
\begin{equation}\label{eq:3-134}
g\alpha\beta_c \to 21{\cdot}911\!\left(\frac{\Omega}{d}\right)^{\!4/3}\!\kappa^{-1} \quad (\Omega\to\infty;\ \text{or}\ \nu\to 0);
\end{equation}
this should be contrasted with the formula,
\begin{equation}\label{eq:3-135}
g\alpha\beta_c = \text{constant}\ \kappa d^{-4},
\end{equation}
valid in the absence of rotation. The dependence of $\beta_c$ on an inverse power of $\nu$, for $\Omega$ fixed and $\nu\to 0$, implies that an \emph{inviscid, ideal fluid, in rotation, is stable with respect to the onset of stationary convection for all adverse temperature gradients}. This is clearly a consequence of the Taylor--Proudman theorem.

\subsection*{(b) The solution for two rigid boundaries}

We shall obtain the solution for this case by the variational method. In view of the symmetry of this problem with respect to the bounding planes, we shall find it convenient to translate the origin of $z$ to be midway between the two planes. The fluid will then be confined between $z = \pm\frac{1}{2}$ and we shall have to seek solutions of equations~\eqref{eq:3-120}--\eqref{eq:3-123} which satisfy the boundary conditions
\begin{equation}\label{eq:3-136}
F = W = DW = Z = 0 \quad \text{for } z = \pm\tfrac{1}{2}.
\end{equation}

It is apparent from the equations and the boundary conditions that the proper solutions of the problem fall into two non-combining groups; these consist of solutions which are even in $W$ and odd in $Z$ and solutions which are odd in $W$ and even in $Z$; we shall call these the even and the odd solutions, respectively. It is also clear that the lowest characteristic value of $R$ will occur among the even solutions; we shall accordingly consider such solutions.

Since $F$ is assumed to be even and is required to vanish at $z = \pm\frac{1}{2}$, we can expand it in a cosine series in the form
\begin{equation}\label{eq:3-137}
F = \sum_m A_m\cos[(2m+1)\pi z],
\end{equation}
where the summation over $m$ may be assumed to run from zero to infinity; but this is not necessary since we shall consider the coefficients $A_m$ as variational parameters.

With the chosen form for $F$, the equation to be solved for $W$ is (cf.\ equation~\eqref{eq:3-123})
\begin{equation}\label{eq:3-138}
[(D^2 - a^2)^3 + \mathscr{T}\,D^2]\,W = -\sum_m A_m\, c_{m+1}\cos[(2m+1)\pi z],
\end{equation}
where
\begin{equation}\label{eq:3-139}
c_{m+1} = (2m+1)^2\pi^2 + a^2.
\end{equation}

In view of the linearity of equation~\eqref{eq:3-138}, we can express $W$ and $Z$ as sums of the form
\begin{equation}\label{eq:3-140}
W = \sum_m A_m\, W_m \quad \text{and} \quad Z = \sum_m A_m\, Z_m,
\end{equation}
where $W_m$ and $Z_m$ are solutions of the equations
\begin{equation}\label{eq:3-141}
(D^2 - a^2)^3 W_m + \mathscr{T}\,D^2 W_m = -c_{m+1}\cos[(2m+1)\pi z]
\end{equation}
and
\begin{equation}\label{eq:3-142}
(D^2 - a^2)Z_m = -\left(\frac{2\Omega}{\nu}\,d\right)\!DW_m.
\end{equation}

The variational principle formulated in \S\,26\,(a) is equivalent to minimising
\begin{equation}\label{eq:3-143}
\int_{-1/2}^{+1/2} [(DF)^2 + a^2 F^2]\,dz
\end{equation}
for such variations of the parameters (the $A_m$'s in the present connexion) which preserve the constancy of
\begin{equation}\label{eq:3-144}
\int_{-1/2}^{+1/2} WF\,dz.
\end{equation}

By arguments similar to those used in Chapter~II, \S\,17, it can be shown, by the use of a Lagrangian undetermined multiplier, that the application of this principle is exactly equivalent to solving a secular determinant which will arise from a Fourier analysis of the equation
\begin{equation}\label{eq:3-145}
(D^2 - a^2)\sum_m A_m\cos[(2m+1)\pi z] = -Ra^2\sum_m A_m\, W_m.
\end{equation}

Returning to equation~\eqref{eq:3-141}, we observe that the general solution of this equation, which is even, can be expressed in the form
\begin{equation}\label{eq:3-146}
W_m = c_{m+1}\,\gamma_{2m+1}\cos[(2m+1)\pi z] + \sum_{j=1}^{3} B_j^{(m)}\cosh q_j\, z,
\end{equation}
where
\begin{equation}\label{eq:3-147}
\frac{1}{\gamma_{2m+1}} = [(2m+1)^2\pi^2 + a^2]^3 + (2m+1)^2\pi^2\,\mathscr{T} = c_{2m+1}^3 + (2m+1)^2\pi^2\,\mathscr{T},
\end{equation}
the $B_j^{(m)}$'s ($j = 1,2,3$) are constants of integration, and the $q_j$'s ($j = 1,2,3$) are the three roots of the cubic equation
\begin{equation}\label{eq:3-148}
(q^2 - a^2)^3 + \mathscr{T}\,q^2 = 0.
\end{equation}

An identity which readily follows from equation~\eqref{eq:3-146} is
\begin{equation}\label{eq:3-149}
\frac{1}{\gamma_{2m+1}} = \frac{1}{\Delta}[(2m+1)^2\pi^2 + q_j^2];
\end{equation}
for,
\begin{equation}\label{eq:3-150}
(q^2 - a^2)^3 + \mathscr{T}\,q^2 = \prod_{j=1}^{3}(q^2 - q_j^2)
\end{equation}
is, in reality, an identity.

The solution for $Z_m$ corresponding to the solution~\eqref{eq:3-146} for $W_m$ is
\begin{align}
Z_m &= -\left(\frac{2\Omega}{\nu}\,d\right)\!\bigl\{(2m+1)\pi\gamma_{2m+1}\sin[(2m+1)\pi z] + \notag\\
&\qquad\qquad + \sum_{j=1}^{3} B_j^{(m)}\frac{q_j}{x_j}\sinh q_j\, z\bigr\}, \label{eq:3-151}
\end{align}
where
\begin{equation}\label{eq:3-152}
x_j = q_j^2 - a^2.
\end{equation}

The boundary conditions $W = DW = Z = 0$ for $z = \pm\frac{1}{2}$ require
\begin{align}
&\sum_{j=1}^{3} B_j^{(m)}\cosh\tfrac{1}{2}q_j = 0, \notag\\[4pt]
&\sum_{j=1}^{3} B_j^{(m)}\,q_j\sinh\tfrac{1}{2}q_j = (-1)^m(2m+1)\pi\gamma_{2m+1}\, c_{m+1}, \label{eq:3-153}\\[4pt]
&\sum_{j=1}^{3} B_j^{(m)}\frac{q_j}{x_j}\sinh\tfrac{1}{2}q_j = (-1)^{m+1}(2m+1)\pi\gamma_{2m+1}. \notag
\end{align}

On solving these equations, we find that
\begin{align}
B_1^{(m)} &= (-1)^m(2m+1)\pi\gamma_{2m+1}\bigl\{q_1\bigl(1+\tfrac{c_{m+1}}{x_1}\bigr)\coth\tfrac{1}{2}q_1\, s - \notag\\
&\qquad\qquad -q_2\bigl(1+\tfrac{c_{m+1}}{x_2}\bigr)\coth\tfrac{1}{2}q_2\,\Delta\operatorname{cosech}\tfrac{1}{2}q_1\bigr\}, \label{eq:3-154}
\end{align}
where
\begin{equation}\label{eq:3-155}
\frac{1}{\Delta} = \frac{q_1 q_2}{x_1 x_2}(x_2 - x_1)\coth\tfrac{1}{2}q_1 + \frac{q_1 q_2}{x_1 x_2}(x_1 - x_2)\coth\tfrac{1}{2}q_2 + \frac{q_1 q_2}{x_1 x_2}(x_2 - x_1)\coth\tfrac{1}{2}q_1,
\end{equation}
and $B_2^{(m)}$ and $B_3^{(m)}$ are given by similar expressions which can be obtained by cyclically permuting the $q_j$'s and the $x_j$'s in the solution~\eqref{eq:3-154}.

The characteristic equation is now obtained from (cf.\ equations~\eqref{eq:3-145} and~\eqref{eq:3-146})
\begin{equation}\label{eq:3-156}
\sum_m A_m\, c_{m+1}\cos[(2m+1)\pi z] = Ra^2\sum_m A_m\!\left[c_{m+1}\,\gamma_{2m+1}\cos[(2m+1)\pi z] + \sum_{j=1}^{3} B_j^{(m)}\cosh q_j\, z\right].
\end{equation}

Multiplying this equation by $\cos[(2n+1)\pi z]$ and integrating over the range of $z$, we obtain
\begin{equation}\label{eq:3-157}
\tfrac{1}{2}c_{m+1}\, A_m = Ra^2\!\left[\tfrac{1}{2}c_{m+1}\,\gamma_{2m+1}\, A_m + \sum_n (n|m)\,A_n\right] \quad (n = 0,1,2,\ldots),
\end{equation}
where
\begin{equation}\label{eq:3-158}
(n|m) = \sum_{j=1}^{3} B_j^{(m)} \int_{-1/2}^{+1/2} \cosh q_j\, z\cos[(2n+1)\pi z]\,dz
= 4m\pi(-1)^n\sum_{j=1}^{3} \frac{B_j^{(m)}\cosh\tfrac{1}{2}q_j}{(2n+1)^2\pi^2 + q_j^2}.
\end{equation}

Equation~\eqref{eq:3-158} provides a set of linear homogeneous equations for the constants $A_m$; this same set of equations ensures that expression~\eqref{eq:3-143} is a minimum for all variations of the $A_m$'s which preserve the constancy of~\eqref{eq:3-144}. In this latter formulation $Ra^2$ is the Lagrangian undetermined multiplier.

The determinant of the system of equations represented by~\eqref{eq:3-157} must vanish; and this provides the characteristic equation for $R$. Thus,
\begin{equation}\label{eq:3-159}
\left\|\tfrac{1}{2}c_{m+1}\!\left[\frac{1}{Ra^2} - \gamma_{2m+1}\right]\!\delta_{mn} - (n|m)\right\| = 0.
\end{equation}

Substituting for the $B_j^{(m)}$'s in accordance with equation~\eqref{eq:3-154} and making use of the identity~\eqref{eq:3-149}, we find after some minor reductions that the expression for $(n|m)$ given in~\eqref{eq:3-158} becomes
\begin{align}
(n|m) &= 2(-1)^{m+n}(2n+1)(2m+1)\pi^2\gamma_{2m+1}\,\gamma_{2n+1}\,\Delta\prod_{j=1}^{3}\coth\tfrac{1}{2}q_j\times \notag\\
&\quad\times\Bigl\{q_1(q_2^2 - q_3^2)\bigl[(2n+1)^2\pi^2 + q_1^2\bigr]\bigl(1+\tfrac{c_{m+1}}{x_1}\bigr)\tanh\tfrac{1}{2}q_1 + \notag\\
&\qquad + q_2(q_1^2 - q_3^2)(2n+1)^2\pi^2 + q_2^2\bigr]\bigl(1+\tfrac{c_{m+1}}{x_2}\bigr)\tanh\tfrac{1}{2}q_2 + \notag\\
&\qquad + q_3(q_1^2 - q_2^2)\bigl[(2n+1)^2\pi^2 + q_3^2\bigr]\bigl(1+\tfrac{c_{m+1}}{x_3}\bigr)\tanh\tfrac{1}{2}q_3\Bigr\}. \label{eq:3-160}
\end{align}

Recalling the definitions of $c_{m+1}$ and $x_j$ (equations~\eqref{eq:3-139} and~\eqref{eq:3-152}), we can rewrite the foregoing in the form
\begin{align}
(n|m) &= 2(-1)^{m+n}(2n+1)(2m+1)\pi^2\gamma_{2m+1}\,\gamma_{2n+1}\,\Delta\prod_{j=1}^{3}\coth\tfrac{1}{2}q_j\times \notag\\
&\quad\times\Bigl\{(x_2 - x_3)(c_{m+1} + x_1)\bigl(1+\tfrac{c_{m+1}}{x_1}\bigr)\,q_1\tanh\tfrac{1}{2}q_1 + \notag\\
&\qquad + (x_1 - x_3)(c_{m+1} + x_2)\bigl(1+\tfrac{c_{m+1}}{x_2}\bigr)\,q_2\tanh\tfrac{1}{2}q_2 + \notag\\
&\qquad + (x_1 - x_2)(c_{m+1} + x_3)\bigl(1+\tfrac{c_{m+1}}{x_3}\bigr)\,q_3\tanh\tfrac{1}{2}q_3\Bigr\}. \label{eq:3-161}
\end{align}

This expression for the matrix element $(n|m)$ is manifestly symmetric in $n$ and $m$: a fact which reflects the basic self-adjoint character of the underlying problem.

The expression for $(n|m)$ can be further reduced. First, we observe that according to equations~\eqref{eq:3-148} and~\eqref{eq:3-152} the $x_j$'s are the roots of the equation
\begin{equation}\label{eq:3-162}
x^3 + \mathscr{T}\,x + \mathscr{T}\,a^2 = 0.
\end{equation}

This equation allows one real root and a pair of complex conjugate roots; let these be
\begin{equation}\label{eq:3-163}
x \quad \text{and} \quad X \pm iY.
\end{equation}

Denote the corresponding roots of $q$ by
\begin{equation}\label{eq:3-164}
q = \sqrt{a^2 + x} \quad \text{and} \quad \alpha_1 \pm i\alpha_2 = \sqrt{a^2 + X \pm iY}.
\end{equation}

On substituting for the $x_j$'s and the $q_j$'s in accordance with these definitions, we find that the expression~\eqref{eq:3-161} for $(n|m)$ can be reduced to the form
\begin{align}
(n|m) &= 2(-1)^{m+1}(2n+1)(2m+1)\pi^2\gamma_{2m+1}\,\gamma_{2n+1}\,\Phi^{(0)}\times \notag\\
&\quad\times\Bigl\{Y(c_{2m+1}+x)\bigl(1+\tfrac{c_{2m+1}}{x}\bigr)\tanh\tfrac{1}{2}q + \notag\\
&\qquad + \tfrac{c_{2n+1}\,c_{2m+1}}{X^2+Y^2}\,\Phi_1^{(c)} - (c_{2n+1}+c_{2m+1})\,\Phi_2^{(c)} + \Phi_3^{(c)}\Bigr\}, \label{eq:3-165}
\end{align}
where $\Phi^{(0)}, \ldots, \Phi_3^{(c)}$ are functions of $q$, $\alpha_1$, $\alpha_2$, $x$, $X$, and $Y$ defined as follows.

Let
\begin{align}
\phi_1^{(c)} &= \frac{\alpha_1\sinh\alpha_1 - \alpha_2\sin\alpha_2}{\cosh\alpha_1 + \cos\alpha_2}, \qquad
\psi_1^{(c)} = \frac{\alpha_2\sinh\alpha_1 + \alpha_1\sin\alpha_2}{\cosh\alpha_1 + \cos\alpha_2}, \notag\\[6pt]
\phi_2^{(c)} &= \frac{\alpha_1\sinh\alpha_1 + \alpha_2\sin\alpha_2}{\cosh\alpha_1 - \cos\alpha_2}, \qquad
\psi_2^{(c)} = \frac{\alpha_2\sinh\alpha_1 - \alpha_1\sin\alpha_2}{\cosh\alpha_1 - \cos\alpha_2}, \label{eq:3-166}
\end{align}
then
\begin{align}
\Phi^{(c)} &= \coth\tfrac{1}{2}q\,(\sinh^2\!\tfrac{1}{2}\alpha_1 + \sin^2\!\tfrac{1}{2}\alpha_2)^{-1}\times \notag\\
&\quad\times\left\{q\frac{Y\phi_1^{(c)} - X\psi_1^{(c)}}{X^2+Y^2} + \frac{\phi_1^{(c)}}{z}\right\} - \frac{\alpha_1+\alpha_2}{X^2+Y^2}\,Y\coth\tfrac{1}{2}q, \label{eq:3-167}
\end{align}
\begin{align}
\Phi_1^{(c)} &= (X^2 - Y^2 - xX)\psi_1^{(c)} - Y(2X-x)\phi_1^{(c)}, \notag\\[4pt]
\Phi_2^{(c)} &= Y\phi_1^{(c)} - (X-x)\psi_1^{(c)}, \notag\\[4pt]
\Phi_3^{(c)} &= (X^2 + Y^2 - xX)\psi_1^{(c)} - xY\phi_1^{(c)}. \label{eq:3-168}
\end{align}

By solving the determinantal equation~\eqref{eq:3-159} by including successively more rows and columns, we can evaluate the required characteristic values of $R$ with increasing precision. The results of such calculations are summarized in Table~VIII and illustrated in Figs.~21 and~22 (see p.~96). It will be seen that in no case is it really necessary to go higher than the second approximation.

\begin{table}[ht]
\centering
\caption{Critical Rayleigh numbers and related constants for the case when both bounding surfaces are rigid and the onset of instability is as stationary convection}
\label{tab:VIII}
\begin{tabular}{cccccccc}
\hline
& & \multicolumn{3}{c}{$R_c$} & \multicolumn{1}{c}{Second} & \multicolumn{2}{c}{Third approximation} \\
\cline{3-5}\cline{7-8}
$\mathscr{T}$ & $a_c$ & First & Second & Third & approxi- & & \\
& & approx. & approx. & approx. & mation $A_1/A_0$ & $A_1/A_0$ & $A_2/A_1$ \\
\hline
$10$ & $3{\cdot}10$ & $1{\cdot}720\times10^3$ & & $+0{\cdot}00881$ & & \\
$100$ & $3{\cdot}15$ & $1{\cdot}764\times10^3$ & & $+0{\cdot}00845$ & & \\
$500$ & $3{\cdot}45$ & $1{\cdot}948\times10^3$ & $1{\cdot}9405\times10^3$ & $+0{\cdot}00840$ & & \\
$1{,}000$ & $3{\cdot}65$ & $1{\cdot}945\times10^3$ & & $+0{\cdot}02349$ & $+0{\cdot}00091$ \\
$3{,}000$ & $3{\cdot}95$ & & & & \\
$5{,}000$ & $4{\cdot}80$ & & $3{\cdot}4688\times10^4$ & $+0{\cdot}00233$ & & \\
$10{,}000$ & & & & & & \\
$10^5$ & & & & $-0{\cdot}00403$ & & \\
$10^6$ & $10{\cdot}65$ & $1{\cdot}737\times10^5$ & & $-0{\cdot}00493$ & & \\
$10^7$ & $14{\cdot}5$ & $3{\cdot}513\times10^5$ & & $+0{\cdot}00055$ & & \\
$10^8$ & $44{\cdot}5$ & $3{\cdot}4698\times10^6$ & $3{\cdot}4574\times10^6$ & $-0{\cdot}00731$ & $-0{\cdot}07798$ & $+0{\cdot}04193$ \\
\hline
\end{tabular}
\end{table}

\subsection*{(c) The solution for one rigid and one free boundary}

In this case the conditions to be satisfied on the two bounding surfaces are different. However, the solution for this case can be reduced to that of case~(b) by considering odd (instead of even) solutions for $W$. For it is clear that an odd solution for $W$ satisfying the boundary conditions appropriate to case~(b) vanishes at $z = 0$; therefore, at $z = 0$ the boundary conditions on $W$ appropriate to a free surface are satisfied. Associated with an odd solution for $W$ is an even solution for $Z$; therefore, at $z = 0$, $DZ$ will vanish which is the further condition to be satisfied on a free surface. Consequently, a solution with $W$ odd and $Z$ even, suitable for case~(b) and applicable to a cell depth~$\frac{1}{2}d$, provides a solution for case~(c) applicable to a cell depth~$d$, and Rayleigh and Taylor numbers which are sixteen times smaller.

We consider then the odd solutions appropriate to case~(b) and obtain them once again by the variational method.

We now expand $F$ in a sine series of the form
\begin{equation}\label{eq:3-168b}
F = \sum_m A_m\sin 2m\pi z
\end{equation}
and follow the procedure outlined in \S\,(b) above. We find that the corresponding solutions for $W_m$ and $z_m$ are
\begin{equation}\label{eq:3-169}
W_m = c_{2m}\,\gamma_{2m}\sin 2m\pi z + \sum_{j=1}^{3} B_j^{(m)}\sinh q_j\, z
\end{equation}
and
\begin{equation}\label{eq:3-170}
Z_m = -\left(\frac{2\Omega}{\nu}\,d\right)\!\left(-2m\pi\gamma_{2m}\cos 2m\pi z + \sum_{j=1}^{3} B_j^{(m)}\frac{q_j}{x_j}\cosh q_j\, z\right),
\end{equation}
where $q_j$ and $x_j$ have the same meanings as in \S\,(b) and $c_{2m}$ and $\gamma_{2m}$ are defined as in equations~\eqref{eq:3-139} and~\eqref{eq:3-147} with $2m$ replacing $(2m+1)$. The constants of integration $B_j^{(m)}$ are similarly determined by the boundary conditions.

Substituting $F$ and $W$ in accordance with equations~\eqref{eq:3-168b} and~\eqref{eq:3-169} in the equation relating them (namely, equation~\eqref{eq:3-105}), we now have (cf.\ equation~\eqref{eq:3-156})
\begin{equation}\label{eq:3-171}
\sum_m A_m\, c_{2m}\sin 2m\pi z = Ra^2\sum_m A_m\!\left[c_{2m}\,\gamma_{2m}\sin 2m\pi z + \sum_{j=1}^{3} B_j^{(m)}\sinh q_j\, z\right].
\end{equation}

Multiplying this equation by $\sin 2n\pi z$ and integrating over the range of $z$, we obtain
\begin{equation}\label{eq:3-172}
\tfrac{1}{2}c_{2m}\, A_m = Ra^2\!\left[\tfrac{1}{2}c_{2m}\,\gamma_{2m}\, A_m + \sum_n (n|m)\,A_n\right],
\end{equation}
where
\begin{align}
(n|m) &= \sum_{j=1}^{3} B_j^{(m)} \int_{-1/2}^{+1/2} \sinh q_j\, z\sin 2m\pi z\,dz \notag\\
&= 4m\pi(-1)^n \sum_{j=1}^{3} \frac{B_j^{(m)}\sinh\tfrac{1}{2}q_j}{4n^2\pi^2 + q_j^2}, \label{eq:3-173}
\end{align}
and this leads to the characteristic equation
\begin{equation}\label{eq:3-174}
\left\|\tfrac{1}{2}c_{2m}\!\left[\frac{1}{Ra^2} - \gamma_{2m}\right]\!\delta_{mn} - (n|m)\right\| = 0.
\end{equation}

The explicit expression for the matrix element $(n|m)$ is found to be (cf.\ equation~\eqref{eq:3-165})
\begin{align}
(n|m) &= 8(-1)^{m+n+1}nm\pi^2\gamma_{2m}\,\gamma_{2n}\,\Phi^{(0)}\times \notag\\
&\quad\times\Bigl\{Y(c_{2m}+x)\bigl(1+\tfrac{c_{2m}}{x}\bigr)\coth\tfrac{1}{2}q + \notag\\
&\qquad + \tfrac{c_{2n}\,c_{2m}}{X^2+Y^2}\,\Phi^{(c)} - (c_{2n}+c_{2m})\,\Phi^{(c)} + \Phi^{(c)}\Bigr\}, \label{eq:3-175}
\end{align}
where $x$, $q$, $X$, $Y$ have the same meanings as in \S\,(b); and analogous to equations~\eqref{eq:3-166} and~\eqref{eq:3-167}, we now have
\begin{align}
\phi_1^{(c)} &= \frac{\alpha_2\sinh\alpha_1 + \alpha_2\sin\alpha_2}{\cosh\alpha_1 + \cos\alpha_2}, \qquad
\psi_1^{(c)} = \frac{\alpha_2\sinh\alpha_1 - \alpha_1\sin\alpha_2}{\cosh\alpha_1 + \cos\alpha_2}, \notag\\[6pt]
\phi_2^{(c)} &= \frac{\alpha_1\sinh\alpha_1 - \alpha_2\sin\alpha_2}{\cosh\alpha_1 - \cos\alpha_2}, \qquad
\psi_2^{(c)} = \frac{\alpha_2\sinh\alpha_1 + \alpha_1\sin\alpha_2}{\cosh\alpha_1 - \cos\alpha_2}, \label{eq:3-176}
\end{align}
\begin{align}
\Phi^{(c)} &= \frac{\tanh\tfrac{1}{2}q\,(\sinh^2\!\tfrac{1}{2}\alpha_1 + \sin^2\!\tfrac{1}{2}\alpha_2)(\cosh\alpha_1 + \cos\alpha_2)^{-1}}{\left\{q\frac{Y\phi^{(c)} - X\psi^{(c)}}{X^2+Y^2} + \frac{\psi^{(c)}}{z}\right\}} - \frac{\alpha_1+\alpha_2}{X^2+Y^2}\,Y\tanh\tfrac{1}{2}q. \notag
\end{align}

\begin{align}
\Phi_1^{(c)} &= (X^2 - Y^2 - xX)\psi_1^{(c)} - Y(2X-x)\phi_1^{(c)}, \notag\\
\Phi_2^{(c)} &= Y\phi_1^{(c)} - (X-x)\psi_1^{(c)}, \notag\\
\Phi_3^{(c)} &= (X^2 + Y^2 - xX)\psi_1^{(c)} - xY\phi_1^{(c)}. \label{eq:3-176b}
\end{align}

The secular equation~\eqref{eq:3-174} has been solved in various approximations. The results of such calculations are summarized in Table~IX and illustrated in Figs.~21 and~22 (see p.~96).

\begin{table}[ht]
\centering
\caption{Critical Rayleigh numbers and related constants for the case when one bounding surface is rigid and the other is free and the onset of instability is as stationary convection}
\label{tab:IX}
\begin{tabular}{cccccccc}
\hline
& & \multicolumn{3}{c}{$R_c$} & \multicolumn{1}{c}{Second} & \multicolumn{2}{c}{Third approximation} \\
\cline{3-5}\cline{7-8}
$\mathscr{T}$ & $a_c$ & First & Second & Third & approxi- & & \\
& & approx. & approx. & approx. & mation $A_1/A_0$ & $A_1/A_0$ & $A_2/A_1$ \\
\hline
$6{\cdot}25$ & $2{\cdot}68$ & $1{\cdot}120\times10^3$ & $1{\cdot}1085\times10^3$ & $+0{\cdot}0680$ & & \\
$3{\cdot}125\times10^1$ & $2{\cdot}70$ & $1{\cdot}148\times10^3$ & $1{\cdot}1395\times10^3$ & $+0{\cdot}06140$ & $+0{\cdot}06335$ & $-0{\cdot}00101$ \\
$1{\cdot}875\times10^2$ & $2{\cdot}975$ & $1{\cdot}393\times10^3$ & & $+0{\cdot}04935$ & & \\
$6{\cdot}250\times10^2$ & $3{\cdot}40$ & $2{\cdot}060\times10^3$ & $1{\cdot}6376\times10^4$ & $+0{\cdot}03570$ & $+0{\cdot}01559$ & $-0{\cdot}01198$ \\
$1{\cdot}875\times10^3$ & $4{\cdot}125$ & $3{\cdot}950\times10^3$ & & $+0{\cdot}01033$ & & \\
$6{\cdot}250\times10^3$ & $6{\cdot}005$ & $7{\cdot}291\times10^3$ & & $+0{\cdot}00604$ & & \\
$6{\cdot}250\times10^4$ & $7{\cdot}475$ & $1{\cdot}453\times10^4$ & $1{\cdot}4510\times10^5$ & $+0{\cdot}03702$ & $-0{\cdot}03701$ & $+0{\cdot}00491$ \\
$1{\cdot}875\times10^5$ & $9{\cdot}00$ & $2{\cdot}850\times10^4$ & & $+0{\cdot}01796$ & & \\
$6{\cdot}250\times10^5$ & $11{\cdot}05$ & $6{\cdot}177\times10^4$ & & $-0{\cdot}02774$ & & \\
& $13{\cdot}80$ & $1{\cdot}381\times10^5$ & $8{\cdot}2265\times10^5$ & $-0{\cdot}07634$ & $-0{\cdot}07907$ & $+0{\cdot}03261$ \\
$10^7$ & $26{\cdot}55$ & $1{\cdot}730\times10^5$ & $1{\cdot}7173\times10^6$ & $-0{\cdot}08461$ & $-0{\cdot}08504$ & $+0{\cdot}04698$ \\
$10^8$ & $58{\cdot}15$ & $3{\cdot}774\times10^5$ & $3{\cdot}7579\times10^6$ & $-0{\cdot}06739$ & $-0{\cdot}06841$ & $+0{\cdot}03876$ \\
\hline
\end{tabular}
\end{table}

\subsection*{(d) The origin of the $\mathscr{T}^{1/3}$- and the $\mathscr{T}^{2/3}$-laws}

From the results for the three sets of boundary conditions illustrated in Figs.~21 and~22, it is apparent that all three cases exhibit the same general features. This is particularly true of the asymptotic dependence on $\mathscr{T}$ of the critical Rayleigh number and the associated wave number. For the case of two free boundaries, the laws
\begin{equation}\label{eq:3-177}
R_c \to \text{constant}\ \mathscr{T}^{2/3} \quad \text{and} \quad a \to \text{constant}\ \mathscr{T}^{1/6}
\end{equation}
follow directly from the solution of the characteristic value problem. From the results of the calculations for the two other cases, it appears that the same power laws hold for them also though the constants of proportionality seem to depend slightly, but definitely, on the boundary conditions.

By going back to the original differential equations, we shall try to locate the common origin of the $\mathscr{T}^{1/3}$- and the $\mathscr{T}^{2/3}$-laws.

Consider the equation
\begin{equation}\label{eq:3-178}
(D^2 - a^2)^3 W + \mathscr{T}\,D^2 W = -Ra^2 W.
\end{equation}

When $\mathscr{T}\to\infty$, we expect that for $R$ in the neighbourhood of $R_c$, $a$ also will tend to infinity. But we do not expect that the solution for $W$, in this limit, will show any special behaviour which will make any of its higher derivatives become `disproportionately' large. The solution for the case of two free boundaries, as well as the progression of the values of the coefficients $A_m$ in the variational solutions for the other two cases (see Tables~VIII and~IX), support this latter expectation. Accordingly, keeping only the terms in $\mathscr{T}$ and in the highest power of $a$ in equation~\eqref{eq:3-178}, as relevant in the limit $\mathscr{T}\to\infty$, we have
\begin{equation}\label{eq:3-179}
\mathscr{T}\,D^2 W = -(Ra^2 - a^6)W.
\end{equation}

This is an equation of the second order in $W$. Consequently, we cannot satisfy all the boundary conditions of the problem; we can satisfy only two of the six boundary conditions. It would appear that the vanishing of $W$ on the boundaries is one condition which we must satisfy on all

accounts. The solution of equation~\eqref{eq:3-179} which vanishes for $z = 0$ and~1 is
\begin{equation}\label{eq:3-180}
W = A\sin n\pi z,
\end{equation}
where $n$ is an integer. The corresponding lowest characteristic value of $R$ is
\begin{equation}\label{eq:3-181}
R = \frac{1}{a^2}(a^6 + \pi^2\,\mathscr{T}).
\end{equation}

As a function of $a$, $R$ given by this equation attains its minimum when
\begin{equation}\label{eq:3-182}
4a^3 - \frac{2}{a^3}\pi^2\,\mathscr{T} = 0,
\end{equation}
or
\begin{equation}\label{eq:3-183}
a = (\tfrac{1}{2}\pi^2\,\mathscr{T})^{1/6};
\end{equation}
and the associated value of $R$ is
\begin{equation}\label{eq:3-184}
R = 3(\tfrac{1}{2}\pi^2\,\mathscr{T})^{2/3}.
\end{equation}

These results agree with those given in equation~\eqref{eq:3-133}.

In one sense, the foregoing arguments do not go much beyond the discussion for the case of two free boundaries. But it does suggest that the origin of the $\mathscr{T}^{1/3}$- and the $\mathscr{T}^{2/3}$-laws probably lies in the effective lowering of the order of equation~\eqref{eq:3-178} and a corresponding reduction in the number of boundary conditions that need be satisfied as $\mathscr{T}\to\infty$. It is worth pointing out that the preceding discussion is not `fine' enough to account for the differences in the constants of proportionality in the asymptotic relations which the variational calculations strongly suggest.

\section{The motions in the horizontal planes and the cell patterns at the onset of instability as stationary convection}\label{sec:3-28}

We now turn to a description of the motions in the horizontal plane and of the cell patterns which can emerge at the onset of instability as stationary convection.

In terms of the solutions for $w$ and $\zeta$, the components of the velocity in the horizontal plane are given by (equations~II~(110) and~(111))
\begin{equation}\label{eq:3-185}
\mathbf{u} = \frac{1}{a^2}\!\left(\frac{\partial^2 w}{\partial x\,\partial z}+\frac{\partial\zeta}{\partial y}\right), \quad
\mathbf{v} = \frac{1}{a^2}\!\left(\frac{\partial^2 w}{\partial y\,\partial z}-\frac{\partial\zeta}{\partial x}\right).
\end{equation}

To be specific, let
\begin{align}
w &= W(z)\cos a_x\, x\cos a_y\, y \label{eq:3-186}\\
\intertext{and}
\zeta &= Z(z)\cos a_x\, x\cos a_y\, y. \notag
\end{align}

Equations~\eqref{eq:3-185} then give
\begin{align}
\mathbf{u} &= -\frac{1}{a^2}(a_x\,DW\sin a_x\, x\cos a_y\, y + a_y\,Z\cos a_x\, x\sin a_y\, y), \notag\\
\mathbf{v} &= -\frac{1}{a^2}(a_y\,DW\cos a_x\, x\sin a_y\, y - a_x\,Z\sin a_x\, x\cos a_y\, y). \label{eq:3-187}
\end{align}

Since (cf.\ equation~\eqref{eq:3-122})
\begin{equation}\label{eq:3-188}
(D^2 - a^2)Z = -\left(\frac{2\Omega d}{\nu}\right)\!DW,
\end{equation}
we can also write
\begin{align}
\mathbf{u} &= \frac{1}{(2\Omega d)/a^2}\!\frac{1}{\nu}\!\bigl\{(a_y\sin a_x\, x\cos a_y\, y)(D^2 - a^2)Z - \notag\\
&\qquad\qquad - (a_x\cos a_x\, x\sin a_y\, y)\,Z\,\sqrt{\mathscr{T}}\bigr\}, \notag\\
\mathbf{v} &= \frac{1}{(2\Omega d)/a^2}\!\frac{1}{\nu}\!\bigl\{(a_y\cos a_x\, x\sin a_y\, y)(D^2 - a^2)Z + \notag\\
&\qquad\qquad + (a_x\sin a_x\, x\cos a_y\, y)\,Z\,\sqrt{\mathscr{T}}\bigr\}. \label{eq:3-189}
\end{align}

From these expressions for $\mathbf{u}$ and $\mathbf{v}$ we find
\begin{equation}\label{eq:3-190}
u^2 + v^2 = \frac{1}{(2\Omega d)^2}\!\frac{1}{\nu^2}(a_x^2\sin^2\!a_x\, x\cos^2\!a_y\, y + a_y^2\cos^2\!a_x\, x\sin^2\!a_y\, y)\times\bigl\{[(D^2 - a^2)Z]^2 + \mathscr{T}\,Z^2\bigr\}.
\end{equation}

Therefore, the average over a horizontal plane, we have
\begin{equation}\label{eq:3-191}
\langle u^2 + v^2\rangle = \frac{1}{4a^2}\frac{1}{(2\Omega d)^2}\bigl\{[(D^2 - a^2)Z]^2 + \mathscr{T}\,Z^2\bigr\}.
\end{equation}

For the case of two free boundaries, the solutions for $W$ and $Z$ appropriate for the lowest mode can be written as (cf.\ equation~\eqref{eq:3-126})
\begin{equation}\label{eq:3-192}
W = W_0\sin\pi z \quad \text{and} \quad Z = W_0\!\left(\frac{2\Omega d}{\nu}\right)\!\frac{\pi}{\pi^2 + a^2}\cos\pi z.
\end{equation}

A somewhat surprising result emerges when the foregoing solution for $Z$ is used in equation~\eqref{eq:3-191}. We find
\begin{equation}\label{eq:3-193}
\langle u^2 + v^2\rangle = \tfrac{1}{4}W_0^2\,\frac{\pi^2}{a^2(\pi^2 + a^2)}\bigl[(\pi^2 + a^2)^2 + \mathscr{T}\bigr]\cos^2\!\pi z.
\end{equation}

Letting $a^2 = \pi^2 x$ (as in \S\,27\,(a)), equation~\eqref{eq:3-129}), we have
\begin{equation}\label{eq:3-194}
\langle u^2 + v^2\rangle = \tfrac{1}{4}W_0^2\,\frac{(1+x)^2 + \mathscr{T}/\pi^4}{x(1+x)^2}\cos^2\!\pi z.
\end{equation}

On the other hand, the value of $x$ at the onset of instability is related to $\mathscr{T}$ by equation~\eqref{eq:3-131}; using this equation to eliminate $\mathscr{T}$ from~\eqref{eq:3-194}, we find that we are left with
\begin{equation}\label{eq:3-195}
\langle u^2 + v^2\rangle = \tfrac{1}{4}W_0^2\cos^2\!\pi z,
\end{equation}
a result which is independent of $\mathscr{T}$. Thus, \emph{for a layer of fluid confined between two free boundaries, the ratio, of the mean kinetic energies of the motions in the horizontal plane and in the vertical direction at the onset of stationary convection, is independent of rotation}.

For other boundary conditions, the invariance expressed by equation~\eqref{eq:3-195} is asymptotically true for $\mathscr{T}\to\infty$: for, according to the discussion in \S\,27\,(d), the general case tends to that for two free boundaries as $\mathscr{T}\to\infty$.

We now turn to the cell patterns. These have been described and beautifully illustrated by Veronis. The following account largely derives from his work.

\subsection*{(a) Rolls, rectangles, and squares}

The case of rectangular cells (of which rolls and squares are special cases) can be derived from the solutions for $\mathbf{u}$ and $\mathbf{v}$ given in~\eqref{eq:3-189}. For the case of two free boundaries (to which the present discussion on the cell patterns will be limited), we can insert for $Z$ its solution from~\eqref{eq:3-192}. We thus obtain the explicit formulae
\begin{align}
\mathbf{u} &= -\frac{\pi}{a^2}\!\left(a_x\sin a_x\, x\cos a_y\, y + \frac{\sqrt{\mathscr{T}}}{\pi^2 + a^2}\,a_y\cos a_x\, x\sin a_y\, y\right)\!W_0\cos\pi z, \notag\\[6pt]
\mathbf{v} &= -\frac{\pi}{a^2}\!\left(a_y\cos a_x\, x\sin a_y\, y - \frac{\sqrt{\mathscr{T}}}{\pi^2 + a^2}\,a_x\sin a_x\, x\cos a_y\, y\right)\!W_0\cos\pi z. \label{eq:3-196}
\end{align}

The corresponding solution for the vertical velocity is given by
\begin{equation}\label{eq:3-197}
w = W_0\cos a_x\, x\cos a_y\, y\sin\pi z.
\end{equation}

The case of rolls is of particular simplicity. The solution for this case can be obtained by setting $a_x = a$ and $a_y = 0$ in equations~\eqref{eq:3-196} and~\eqref{eq:3-197}; thus
\begin{align}
\mathbf{u} &= -\frac{\pi}{a(\pi^2 + a^2)}W_0\sin ax\cos\pi z, \label{eq:3-198}\\
\mathbf{v} &= +\frac{\sqrt{\mathscr{T}}}{a(\pi^2 + a^2)}W_0\sin ax\cos\pi z, \label{eq:3-199}
\end{align}
and
\begin{equation}
w = W_0\cos ax\sin\pi z. \notag
\end{equation}

When $\mathscr{T} = 0$, there are no motions in the $y$-direction and the streamlines are confined to planes normal to the direction of the rolls. When the system is rotated, the Coriolis acceleration induces longitudinal motions. Since in the present case
\begin{equation}\label{eq:3-200}
\frac{v}{u} = -\frac{\sqrt{\mathscr{T}}}{\pi^2 + a^2} = \text{constant},
\end{equation}
the streamlines are again confined to planes; but these planes are inclined to the $z$-axis (see Fig.~23\,$b,c$). Under these circumstances we can define a wave number $a_s$ which will describe the motions in the oblique planes containing the streamlines; it is given by
\begin{equation}\label{eq:3-201}
a_s = a\cos(\tan^{-1} v/u) = \frac{a}{\sqrt{1+v^2/u^2}}.
\end{equation}

Substituting for $v/u$ from~\eqref{eq:3-200}, we obtain
\begin{equation}\label{eq:3-202}
a_s^2 = \frac{a^2(\pi^2 + a^2)^2}{(\pi^2 + a^2)^2 + \mathscr{T}}.
\end{equation}

\figurePlaceholder{23}{(a) A top view of two-dimensional rolls in a non-rotating fluid. (b) The same view in a system rotating counterclockwise. (c) A perspective view of the particle motions in a roll in the rotating case. The arrows indicate the direction of particle motions.}

The quantity on the right-hand side of this equation is, apart from a numerical factor, the reciprocal of the expression which we have already shown to be independent of $\mathscr{T}$ (cf.\ the arguments leading from equation~\eqref{eq:3-193} to~\eqref{eq:3-195}); it has the value $\frac{1}{4}\pi^2$. Thus,
\begin{equation}\label{eq:3-203}
a_s^2 = \tfrac{1}{4}\pi^2;
\end{equation}
but this is the same as the formula which gives the wave number of the roll in the absence of rotation (see equation~II~(195)). Hence, \emph{the wavelength of the roll measured in the plane containing the streamlines is independent of rotation}. This remarkable result is due to Veronis.

Considering next the case of square cells, we obtain the corresponding solutions by setting $a_x = a_y = a/\sqrt{2}$ in~\eqref{eq:3-196}. We obtain
\begin{equation}\label{eq:3-204}
\mathbf{u} = -\frac{\pi}{a\sqrt{2}}\!\left(\sin\frac{ax}{\sqrt{2}}\cos\frac{ay}{\sqrt{2}} + \frac{\sqrt{\mathscr{T}}}{\pi^2 + a^2}\cos\frac{ax}{\sqrt{2}}\sin\frac{ay}{\sqrt{2}}\right)\!W_0\cos\pi z,
\end{equation}
and similarly for $\mathbf{v}$.

The streamlines in the absence of rotation in such square cells have been

\figurePlaceholder{24a}{A square cell in a rotating fluid. The particle motions are from the centre outward along spiral curves. The dashed curves form the boundary of the square cell.}

\noindent illustrated in Chapter~II (Fig.~6, p.~46). The dashed curves in Fig.~24\,$a$ correspond to the sides of the square in Fig.~6. In the presence of rotation, the Coriolis acceleration causes the fluid to turn as it moves towards or outwards from the centre; motions along the curves of constant $w$, as well as transverse to them, occur. Fig.~24\,$b$ is a perspective of the complete trajectory as sketched by Veronis. The drawing illustrates how a fluid element starting at the centre of a cell spirals upwards in a clockwise direction (for a counter-clockwise rotation $\Omega$), and as it approaches the top it crosses over towards a corner of a cell and starts

\figurePlaceholder{24b}{A perspective sketch of the path of a fluid particle in a square cell.}

\noindent spiralling downwards in a counter-clockwise direction. When it arrives at the mid-plane ($z = \frac{1}{2}$), it reverses its sense of spiralling and proceeds downwards, and as it reaches the bottom it crosses back towards the centre of the cell, and starts spiralling upwards towards the centre in a clockwise direction. The reversal of the rotation at the mid-plane is associated with the circumstance that here $DW$ and, therefore, also, $\nabla_1\cdot\mathbf{u}_1$ change sign.

The square cells shown in Fig.~24\,$b$ correspond to the basic geometry of the vertical velocity. The distortion of the cell consequent on the rotation is not shown.

\subsection*{(b) Hexagons}

Christopherson's solution for a hexagonal pattern and appropriate for two free boundaries is (cf.\ equation~II~(248))
\begin{equation}\label{eq:3-205}
w = \frac{1}{3}\!\left(2\cos\frac{2\pi}{L\sqrt{3}}\,x\cos\frac{2\pi}{3L}\,y + \cos\frac{4\pi}{3L}\,y\right)\!W_0\sin\pi z.
\end{equation}

The corresponding solution for the vertical component of the vorticity is
\begin{equation}\label{eq:3-206}
\zeta = \frac{1}{3}\!\left(\frac{2\Omega d}{\nu}\right)\!\frac{\pi}{\pi^2 + a^2}\!\left(2\cos\frac{2\pi}{L\sqrt{3}}\,x\cos\frac{2\pi}{3L}\,y + \cos\frac{4\pi}{3L}\,y\right)\!W_0\cos\pi z.
\end{equation}

Substituting for $w$ and $\zeta$ in accordance with these solutions in~\eqref{eq:3-187} we obtain
\begin{align}
\mathbf{u} &= -\frac{\pi}{3a^2}\!\left(\frac{4\pi}{L\sqrt{3}}\sin\frac{2\pi}{L\sqrt{3}}\,x\cos\frac{2\pi}{3L}\,y + \right.\notag\\
&\qquad\left. +\frac{4\pi}{3L}\frac{\sqrt{\mathscr{T}}}{\pi^2 + a^2}\!\left(\cos\frac{2\pi}{L\sqrt{3}}\,x + 2\cos\frac{2\pi}{3L}\,y\right)\!\sin\frac{2\pi}{3L}\,y\right)\!W_0\cos\pi z, \notag\\[8pt]
\mathbf{v} &= -\frac{\pi}{3a^2}\!\left(\frac{4\pi}{3L}\!\left(\cos\frac{2\pi}{L\sqrt{3}}\,x + 2\cos\frac{2\pi}{3L}\,y\right)\!\sin\frac{2\pi}{3L}\,y - \right.\notag\\
&\qquad\left. -\frac{4\pi}{L\sqrt{3}}\frac{\sqrt{\mathscr{T}}}{\pi^2 + a^2}\sin\frac{2\pi}{L\sqrt{3}}\,x\cos\frac{2\pi}{3L}\,y\right)\!W_0\cos\pi z. \label{eq:3-207}
\end{align}

\figurePlaceholder{25a}{A top view of seven rotating hexagonal cells. The fluid particles follow spiral paths from the centre toward the corners. The dashed lines form the boundaries of the centre cell.}

Fig.~25\,$a$ shows the streamlines as they would appear at the top. This pattern should be contrasted with that illustrated in Fig.~7\,$b$ (p.~50) for a non-rotating system. The dashed spirals form the boundary of the central cell. It will be observed that the source at the centre of the cell provides for the flow downwards at the six corners; and each corner

\figurePlaceholder{25b}{A perspective sketch of the path of a fluid particle in a hexagonal cell.}

\noindent receives fluid from three adjacent cells. In this last respect it is not different from the non-rotating case.

A perspective drawing of the complete trajectory is shown in Fig.~25\,$b$. And finally, in Fig.~26 we illustrate the manner in which the cell is distorted by the rotation.

\figurePlaceholder{26}{A perspective sketch of a hexagonal cell as it is distorted by the rotation of the fluid.}

\subsection*{(c) The limiting behaviour of the streamlines for $\mathscr{T}\to\infty$}

It is clear from the foregoing illustrations of the cell patterns that as $\mathscr{T}$ increases, the streamlines become increasingly close wound spirals. In the limit $\mathscr{T}\to\infty$, when convection in fact ceases, the streamlines become closed curves. Thus, for $\mathscr{T}$ sufficiently large, the dominant terms in the solution~\eqref{eq:3-196} for $\mathbf{u}$ and $\mathbf{v}$ are, for example,
\begin{align}
\mathbf{u} &\to -\frac{\pi\sqrt{\mathscr{T}}}{a^2(\pi^2+a^2)}\,a_y\,W_0\cos a_x\, x\sin a_y\, y\cos\pi z, \notag\\[4pt]
\mathbf{v} &\to +\frac{\pi\sqrt{\mathscr{T}}}{a^2(\pi^2+a^2)}\,a_x\,W_0\sin a_x\, x\cos a_y\, y\cos\pi z. \label{eq:3-208}
\end{align}

The corresponding equation for the streamlines is
\begin{equation}\label{eq:3-209}
\cos a_x\, x\cos a_y\, y = \text{constant}
\end{equation}
or,
\begin{equation}\label{eq:3-210}
w = \text{constant} \quad (\text{for } z = \text{constant}).
\end{equation}

This last result is not restricted to rectangular cells: it applies equally to all cell patterns.

Whereas, in the absence of rotation $\mathbf{u}_\perp$ is parallel to $\nabla_1 w$, in the limit $\mathscr{T}\to\infty$ the streamlines tend to coincide with the curves $w = \text{constant}$. For $\mathscr{T}$ large but finite, we may describe the motions in the horizontal plane as consisting, principally, of motions around the curves of constant $w$; and superposed on this is a slow outward, or inward, directed radial motion which induces a spiral pattern. Near the centres of the cells, the spiral patterns tend to become equiangular. It is the small radial motion in the direction of $\nabla_1 w$ that is responsible for all phenomena associated with convection when $\mathscr{T}$ becomes large.
